% !TEX program = lualatex
% =============================================================================
%  Phase Transitions: Metal--Insulator transition in VO2 and
%  Ferroelectric transition in BaTiO3
%  Advanced Laboratory, Dept. of Physics & Astronomy, Stony Brook University
%
%  Accessible LaTeX edition: main.tex is untagged, main-tagged.tex is tagged PDF/UA-2.
%  Compile with lualatex (>= TeX Live 2024).  Run twice for cross-references.
% =============================================================================

% ---- Accessibility / PDF standard -------------------------------------------
% \DocumentMetadata MUST come before \documentclass.  It switches on the new LaTeX
% tagged-PDF engine, declares the document language, and asks for a PDF 2.0 file
% conforming to PDF/UA-2 (the accessibility standard).
%
% This file builds WITHOUT tags: it is faster and the PDF is smaller.
% main-tagged.tex defines \tagged and inputs this file, which turns on the tags
% and the PDF/UA-2 claim.  The lines below only assemble the key list -- nothing
% is typeset or loaded before \DocumentMetadata -- and the list has to be
% expanded first, because \DocumentMetadata does not expand macros inside it.
%
% tagging=off must come AFTER testphase: that key runs tagging=on, so an earlier
% tagging=off is silently overridden.
\def\metadatakeys{lang = en-US, pdfversion = 2.0,
  testphase = {phase-III, math, graphic, table, firstaid}}
\ifdefined\tagged
  \edef\metadatakeys{\metadatakeys, pdfstandard = ua-2}
\else
  \edef\metadatakeys{\metadatakeys, tagging = off}
\fi
\expandafter\DocumentMetadata\expandafter{\metadatakeys}

\documentclass[11pt,a4paper]{article}

% ---- Fonts -------------------------------------------------------------------
% LuaLaTeX reads UTF-8 and loads Latin Modern (OpenType) by itself; unicode-math
% sets the maths in Latin Modern Math, so formulas carry real Unicode characters.
\usepackage{microtype}

% ---- Page geometry -----------------------------------------------------------
\usepackage[margin=2.5cm,headheight=14pt]{geometry}

% ---- Maths -------------------------------------------------------------------
\usepackage{amsmath}
\usepackage{unicode-math}   % after amsmath; replaces amssymb
\setmathfont{Latin Modern Math}
\usepackage{siunitx}
\sisetup{separate-uncertainty, range-units=single}
% Typeset the percent sign as text.  A bare % inside maths ends up unescaped in
% luamml's MathML cache file, and commenting out the rest of that line breaks
% the cache when the next pass reads it back.
\DeclareSIUnit{\percent}{\text{\%}}

% ---- Graphics, color, floats ------------------------------------------------
\usepackage{graphicx}
\graphicspath{{figures/}{tikz/}}
\usepackage{xcolor}
\usepackage{float}
% Pack floats more tightly so that a single figure never claims a whole page.
\renewcommand{\topfraction}{0.9}
\renewcommand{\bottomfraction}{0.8}
\renewcommand{\textfraction}{0.07}
\renewcommand{\floatpagefraction}{0.75}
\setcounter{topnumber}{3}
\setcounter{bottomnumber}{2}
\setcounter{totalnumber}{4}
\usepackage{booktabs}

% ---- Links -------------------------------------------------------------------
\usepackage{hyperref}
\hypersetup{
  colorlinks = true,
  linkcolor  = [rgb]{0.00,0.25,0.60},
  urlcolor   = [rgb]{0.00,0.35,0.30},
  citecolor  = [rgb]{0.00,0.25,0.60},
  pdftitle   = {Phase Transitions: Metal-Insulator Transition in VO2 and
                Ferroelectric Transition in BaTiO3},
  pdfauthor  = {Michael Gurvitch and Laszlo Mihaly},
  pdfsubject = {Advanced Laboratory write-up, Stony Brook University},
  pdfkeywords= {phase transition, metal-insulator transition, vanadium dioxide,
                ferroelectricity, barium titanate, lock-in amplifier,
                four-probe method, Landau theory},
}
\usepackage{url}

% ---- Convenience macros ------------------------------------------------------
\newcommand{\VOtwo}{VO\textsubscript{2}}
\newcommand{\BTO}{BaTiO\textsubscript{3}}
\newcommand{\Rs}{R_S}
\newcommand{\degC}{\si{\degreeCelsius}}

% \tikzfig{<file base>}{<alt text>}  -- include a standalone TikZ figure that
% has been pre-compiled to tikz/<base>.pdf, together with its alternative text.
% The optional argument holds graphicx keys (default: full text width).
\newcommand{\tikzfig}[3][width=\linewidth]{%
  \includegraphics[#1,alt={#3}]{tikz/#2.pdf}}

% \photo{<file>}{<alt text>} -- a raster image plus its alternative text.
\newcommand{\photo}[3][width=0.7\linewidth]{%
  \includegraphics[#1,alt={#3}]{figures/#2}}

% A boxed safety / practical note, announced in words as well as by its frame.
\newcommand{\notebox}[2]{%
  \par\medskip\noindent
  \fbox{\parbox{\dimexpr\linewidth-2\fboxsep-2\fboxrule}{\textbf{#1.} #2}}%
  \par\medskip}

\title{Phase Transitions:\\
       Metal--Insulator Transition in \VOtwo\ and\\
       Ferroelectric Transition in \BTO}
\author{Advanced Laboratory\\
        Department of Physics \& Astronomy\\
        Stony Brook University}
\date{Write-up version of September 2026\\[2pt]
      \small Sections on ferroelectricity in \BTO\ by Laszlo Mihaly and, in
      part, Michael Gurvitch; sections on phase-transition theory and the
      metal--insulator transition in \VOtwo\ by Michael Gurvitch.}

\begin{document}

\maketitle
\thispagestyle{empty}

\clearpage
\tableofcontents
\clearpage

% =============================================================================
\section*{Pre-experiment and mid-experiment evaluation sheet}
\addcontentsline{toc}{section}{Pre-experiment and mid-experiment evaluation sheet}
% =============================================================================

You must meet with a staff member for the experiment you are doing before
beginning an experiment and before the lab period indicated in the middle of the
experiment for an evaluation of your preparation and your progress to date. This
form must be signed and dated by the staff member to indicate you have done the
required work to date. It must be attached as the first page of your report. You
will lose one grade point from your report grade for each section of this form
that is not properly signed on time. If this is your first experiment you should
also be prepared to demonstrate you have read and understood the information on
log books in the course notes.

\bigskip
\noindent Student name\ \hrulefill

\medskip
\noindent Partner name\ \hrulefill

\subsection*{Pre-lab discussion}

Be prepared to demonstrate you have read the background material for this
experiment (most of this material is contained in this write-up, but also in
cited references; it includes theory described in Appendix~\ref{app:ac}).
Discuss the following questions with the staff:

\begin{enumerate}
  \item Describe what we mean by phase transitions; give some examples; what is
        the difference between phase transitions of the first and second kind?
        Experimentally, in this lab, how can we tell which transition it is?
  \item Give a general description of a metal--insulator transition; what
        happens to the electronic spectrum of a solid when this transition takes
        place? What is the most direct manifestation of this transition, and how
        can it be detected? Is this transition taking place all at once in the
        whole sample, or not (explain)?
  \item Explain what ferroelectricity is and what the mechanism for it in \BTO\
        is.
  \item What does a PID controller do? Be specific about its principle of
        operation.
  \item What is the general principle of operation of a lock-in amplifier?
\end{enumerate}

\bigskip
\noindent Completed before first period for the experiment?\qquad Yes\,/\,No

\bigskip
\noindent Staff comments\ \hrulefill

\bigskip
\noindent Staff signature\ \hrulefill\ Date\ \hrulefill

\clearpage
% =============================================================================
\section{How to prepare for the measurement}
% =============================================================================

This write-up contains a lot of material. You should study parts of it before
coming to the lab, and read and study the entire text throughout the course of
this experiment.

\notebox{Note}{A big chunk of this manual is in the Appendices, so do not treat
them as secondary-importance material.}

You should also supplement this write-up by looking up and reading some of the
books/articles/websites from the list of references. We propose that you start
in the following order:

First read Sections \ref{sec:prepare}--\ref{sec:report}, which will give you an
overview of this lab. Then spend some time learning how to use the lock-in
amplifier. Your first task will be to master preliminary practice measurements
of resistance and capacitance. During the first two periods of the lab work,
concentrate on Section~\ref{sec:techniques} of this write-up entitled
``Measurement techniques'' and study Appendices \ref{app:ac} and
\ref{app:fourprobe}. Read the equipment
manual for the lock-in amplifier and Ref.~\cite{ref:lockin}; study the
principle of its operation. Simultaneously, study the instrument hands-on:
connect a circuit and first measure a known resistor, and then a known
capacitor, using the lock-in.

Once you get familiarized with using the lock-in, read about the rest of the
experimental set-up, which is described in Sections
\ref{sec:apparatus}--\ref{sec:advice}. Then you should start reading more of a
physics theory, which is to be found mainly in a large
Appendix~\ref{app:physics}. The physics of phase transitions is not trivial,
and it will take time to ``digest'' and understand it.

From the beginning, keep good records in your notebook. All electrical circuits
that you will build, simple calculations which accompany construction of these
circuits, such as estimates of the measuring current, choice of a limiting
resistor, calculations of resistance and capacitance, sketches of samples or
experimental apparatus, details of sample preparation and geometry, all notes
concerning physics, and all your calculations of physical quantities should be
recorded in the notebook. Short comments you may write should be
understandable, in case you or someone else will later examine your records. The actual experimental
data will be stored in the computer and need not be in the notebook.

% =============================================================================
\section{Goals of the experiment}\label{sec:prepare}
% =============================================================================

\begin{itemize}
  \item Understand experimental techniques used in this lab, including the
        principle and practical operation of a lock-in amplifier when measuring
        resistance and capacitance, the four-probe method, temperature control.
  \item Learn to handle and prepare samples; learn how to quickly check samples
        before they are used in lengthy measurements.
  \item Study the metal--insulator transition in \VOtwo.
  \item Study the ferroelectric transition in \BTO.
  \item Learn basic physics associated with phase transitions and apply it to
        interpreting specific experiments at hand.
  \item We do not just do measurements for the sake of measurements. An
        experiment should lead to some conclusions and it should culminate and
        result in learning of some ``new'' physics (new to us in this lab;
        perhaps not new to the physics community). Therefore it is very
        important to use your data to extract some meaningful physical results:
        namely, estimates of
        microscopic or theoretical quantities which enhance our quantitative
        understanding of the studied phenomena. You will do so for both phase
        transitions studied in this lab.
\end{itemize}

% =============================================================================
\section{Experimental set-up}
% =============================================================================

\begin{itemize}
  \item Sample chamber with four electric leads.
  \item Temperature controller.
  \item Lock-in amplifier to measure the samples.
  \item Keithley four-probe resistance meter to measure temperature.
  \item Computer (username \texttt{gradlab}, password \texttt{senior}).
  \item Decade resistance box, additional resistors, capacitors, BNC cables etc.
\end{itemize}

See the equipment manuals for details. The temperature controller has a web
interface. Use this interface for programming the controller (see
Appendix~\ref{app:tempcontrol}). The lock-in amplifier and the Keithley
multimeter are connected to a computer by the IEEE~488 interface. A data-taking
routine, written in LabVIEW, is provided for recording data
(Appendix~\ref{app:labview}).

\begin{figure}[htbp]
  \centering
  \tikzfig[width=\linewidth]{fig-daq-wiring}{%
    Wiring diagram of the experiment. At the top left a box labeled Lock-in
    amplifier, EG and G 5302, has measuring-input and oscillator-output
    terminals; the oscillator drives the sample through resistor R-zero and the
    measuring input sits across the sample. Below it a dashed box marks the
    sample chamber, containing the sample capacitor, a heater winding and a
    Pt100 probe. The probe runs to the PID controller box, an Omega CND3, which
    lists probe input terminals 10, 11 and 12, output OUT1 on terminals 8 and 9,
    and RS-485 on terminals 13 (D minus) and 14 (D plus). OUT1 switches a
    solid-state relay that powers the heater from AC mains. On the right, the
    lock-in reaches the lab PC over GPIB at address 12 through an NI adapter,
    and the controller reaches it over RS-485 through a Dtech USB converter.}
  \caption{Instrument connections for the temperature-dependent measurement.}
  \label{fig:daqwiring}
\end{figure}

\paragraph{Samples.}
You will be studying samples of \VOtwo\ and \BTO. If two students work on this
lab, they should together measure at least one \VOtwo\ sample, which will be
provided, and prepare themselves and measure two \BTO\ samples. Sample
preparation and all the measurements can be done in collaboration, but for the
purposes of data analysis we will consider that one of these two samples belongs
to one student, and another one to the other, so that each student will report
and evaluate the data on the common \VOtwo\ sample and on their own \BTO\
sample.

% =============================================================================
\section{Principle of the measurement and an overview of data analysis}
% =============================================================================

Details on these subjects will be found in the rest of this write-up; and you
know what they say: the devil is in the detail. Here is the brief summary:

You will be studying two different phase transitions, both taking place as a
function of temperature: a metal--insulator transition in vanadium dioxide
(\VOtwo) and a ferroelectric transition in barium titanate (\BTO). Read
Appendix~\ref{app:physics} where the physics of phase transitions is reviewed
and detailed.

The sample is placed in a temperature-controlled chamber. The temperature is
slowly changed, first rising, and then cooling down, while either resistance
$R$ (in \VOtwo) or capacitance $C$ (with \BTO\ serving as dielectric) is
measured and data is recorded. The temperature will change between room
temperature and about \SI{100}{\degreeCelsius} in the case of \VOtwo, and
between room temperature and about \SI{150}{\degreeCelsius} in the case of
\BTO. The measurement of $R$ and $C$ is performed in an AC mode using a lock-in
amplifier, which is described below in Section~\ref{sec:techniques} and in
Appendix~\ref{app:ac}.

What will you be able to learn from these measurements? In the case of the
resistive \VOtwo\ data, you will be able to make conclusions concerning the
nature and order of a phase transition, including estimates of the
metal--insulator transition temperature, of the percolation transition
temperature, of the coercive temperature. You will also be able to analyze the
data to extract the gap of \VOtwo\ in a semiconducting phase and to
estimate absolute resistivities of the two phases.

In the case of a capacitance measurement, \BTO\ will be the dielectric inside a
capacitor. In a parallel-plate capacitor,
\begin{equation}
  C = \frac{\varepsilon A}{d} = \frac{\varepsilon_0 k A}{d},
\end{equation}
where $A$ is the area of the capacitor, $d$ is the distance between the plates,
$\varepsilon$ is the dielectric constant, $\varepsilon_0$ is a constant called
the permittivity of free space, and $k$ is the relative dielectric constant. We
will be interested in finding $k$ as a function of temperature. It will not be
possible to exactly measure $A$ and $d$ in the type of samples you will have,
but they can be roughly estimated. So, by measuring $C$ (which is measured quite
precisely), you will be able to estimate $k$ to within about a factor of two
accuracy, the uncertainty coming almost entirely from a limited ability here to
evaluate capacitor geometry. At the same time, you can certainly follow how $k$
will change with temperature with the full precision of the capacitance
measurement. This $T$-dependence will pinpoint the temperature of a phase
transition, $T_C$. Combined, these data will allow a calculation of microscopic parameters
of a ferroelectric transition as outlined in Appendix~\ref{app:physics}. At the
end, you will be able to estimate an interesting microscopic quantity: the very
small distance by which atoms shift in the lattice of \BTO\ in a ferroelectric
transition. The ability to estimate such a quantity is what makes this
experiment especially interesting.

\subsection*{List of physical quantities which you will be expected to extract,
calculate and estimate from your data}

\noindent(Some of these can be measured rather precisely; others only roughly
estimated.)

\paragraph{In \VOtwo:}
\begin{itemize}
  \item The semiconducting gap $E_g$
  \item The resistance ratio between semiconducting and metallic phases
  \item Absolute resistivities in semiconducting and metallic phases
  \item The value of $T_C$
  \item The coercive temperature $T^{*}$
  \item The temperature of a percolation transition $T_M$
\end{itemize}

\paragraph{In \BTO:}
\begin{itemize}
  \item $T_C$
  \item Relative dielectric constant $k$ vs.\ $T$, from about
        \SI{20}{\degreeCelsius} to \SI{150}{\degreeCelsius}
  \item Curie--Weiss law at $T > T_C$ and Curie--Weiss constant $C_{CW}$
  \item The value of a microscopic dipole moment $p$
  \item The effective atomic displacement $x$ responsible for the appearance of
        $p$
  \item The last two quantities can be also calculated for \BTO\ single crystals
        using the literature value of $C_{CW}$
\end{itemize}

% =============================================================================
\section{Looking ahead: what should and should not be included in your Final
         Report}\label{sec:report}
% =============================================================================

Your performance will be judged primarily from the contents of your Final
Report. In it there will be several different (but related) parts:

\begin{itemize}
  \item A physics introduction concerning phase transitions and materials
        which you studied and a physics analysis of your results at the end. To
        get these right, you have to study Appendix~\ref{app:physics} of this
        write-up and related literature.

  \item A description of experimental set-up (including electrical circuit
        diagrams) and experimental techniques that you used (lock-in and how it
        was used; four-probe method), including significant details of how you
        obtained your data. For example, if you are measuring some voltage, but
        your ``final product'' is capacitance, you will need to
        explain in some detail how you proceeded from that voltage to knowing
        that capacitance. This may include phasor diagrams and some formula
        derivations given in Section~\ref{sec:techniques} and
        Appendix~\ref{app:ac}.

  \item Initially you will be measuring \emph{known} resistances and
        capacitances just to learn how to connect the circuits and how to use
        the lock-in amplifier. It is very important to perform such calibrations
        or baseline measurements in order to be sure that later you will be
        measuring your samples correctly. While a detailed description of such
        calibrations may be perfectly appropriate for your mid-lab report, it
        should not take center-stage in the Final Report, except as a short
        mention of the fact that these calibrations were done. You may want to
        report the uncertainty (i.e.\ \% of relative deviation) in the
        measurement of a known $R$ and $C$ (of course, assuming that the values
        of these $R$ and $C$ were known very accurately in the first place). But
        detailed data on such baseline measurements would be out of place; the
        focus of the Final Report should be on measuring real samples which
        exhibit phase transitions.

  \item You will present results of sample measurements, your data, in the form
        of graphs, possibly some tables (but do not list all of the data points
        in long tables --- this is neither necessary nor desirable in the Final
        Report). If you are measuring resistance $R$ and capacitance
        $C$ as a function of temperature $T$, your data may consist of graphs of
        $R(T)$ and $C(T)$.

        Let us discuss this in a little more detail. There exists a notion of
        \emph{raw data}, the data you directly obtain in your measurement,
        directly read off the instruments. And then there is data which was to
        some degree analyzed and organized. For example, you may be measuring
        separate values of voltage $V$ and current $I$, and later you may be
        combining the two when obtaining resistance $R = V/I$. But your interest
        --- and the interest of the reader of your report --- is in that
        resistance $R$, and not in $V$ or $I$. Thus one should use common sense
        when reporting data: probably in this example, separate values of $V$
        and $I$ should not be reported (they can remain in your electronic data
        storage, to be examined if necessary), and $R$ should be reported.
        Likewise, if you are measuring a voltage signal which is being converted
        to a capacitance, do not report those voltages; report capacitance
        directly. You may however mention the range of the voltage signals
        that you measured in order to tell your reader something meaningful
        about the quality of your data.

  \item The data should be presented with uncertainties, which not only should
        be included, but accompanied with a brief explanation of how they were
        obtained or estimated. If a calculation is involved, the error must be
        propagated from the raw data to the calculated quantity. If curve
        fitting is involved, the uncertainty will come out of the fitting
        routine. You need also to point out which uncertainties are more
        significant, and which can be ignored; in this lab this last comment
        will be particularly relevant, as some of the uncertainties will far
        exceed others.

  \item There will be data analysis: corrections which you introduced to your
        data, extraction of physical parameters of interest. And at the end
        there will be some conclusions based on your results. You will need to
        study this write-up and references in order to ``connect the dots'', to
        make your experiments meaningful. When applicable, you will compare your
        results to the published results. Finding relevant publications by doing
        literature searches, reading them, and extracting from these
        publications important parameters, is part of a skill of an experimental
        physicist.
\end{itemize}

% =============================================================================
\section{Measurement techniques}\label{sec:techniques}
% =============================================================================

The main instrument you will be using in your measurements in this lab is
called a \emph{lock-in amplifier}~\cite{ref:lockin}. This wonderful instrument
was invented in the 1950's by a physicist Robert Dicke; it is used extensively
for a variety of measurements in physics experiments and elsewhere.

\begin{figure}[htbp]
  \centering
  \photo[width=0.24\linewidth]{dicke.png}{%
    Black-and-white portrait photograph of the physicist Robert Dicke, seated
    and facing the camera.}
  \caption{Robert Dicke (1916--1997), a famous American physicist who made
    important contributions to the fields of astrophysics, atomic physics,
    cosmology and gravity. He also invented the lock-in amplifier in the 1950s,
    with the first commercial model developed in 1962.}
\end{figure}

The advantage a lock-in gives to an experimentalist is in that it rejects a good
deal of noise, increasing the signal-to-noise ratio and allowing the use of
small currents which will not damage even the most delicate of samples.
Additionally, by virtue of always measuring in an AC mode, it automatically
averages the readout voltage, getting rid of the so-called thermal emf offsets
and DC level drifts which are endemic with DC measurements, even when we are
using the best and most accurate DC voltmeters.

While the lock-in amplifier was invented to measure weak signals in the
presence of large noise, in our measurements we will not be pushing the limits
of its capacity in that regard. Instead, we will be using its ability to
measure signals of different specified phases, as explained below and in
Appendix~\ref{app:ac}. This lab gives you an opportunity to learn about this
important instrument.

The name ``lock-in'' refers to locking to a frequency and to a phase of an AC
signal, and measuring only that part of a signal which is in a specified phase
relationship to the reference. This is a version of the \emph{phase-sensitive
detection} technique. The phase can be continually varied and also changed in
\ang{90} steps.

You should read about lock-in operation in our lock-in manual (available online
through Blackboard) and elsewhere~\cite{ref:lockin}. Here and in
Appendix~\ref{app:ac} we give a brief general introduction to its operating
principle, and discuss in more detail the two types of measurements we will be
doing in this lab.

\notebox{Review AC circuits first}{If you have not been dealing with AC
circuits lately, this would be a good time to thoroughly review them: read
about AC currents in general, and the phasor method in particular, in your
introductory physics course textbook; some of this is also discussed in
Ref.~\cite{ref:accircuits}. Study this section and Appendix~\ref{app:ac}; it is
essential that you understand this topic in order to do this lab.}

\paragraph{The principle of lock-in measurement} is as follows: the lock-in
contains its own AC voltage source called \emph{oscillator output}, in which you
can adjust the amplitude (but only up to \SI{2.0}{\volt}), and frequency
(adjustable in a wide range from a fraction of \SI{1}{\hertz} to
\SI{100}{\kilo\hertz}). The frequency, once you set it, will be kept with great
precision and stability by a lock-in. The readout is tuned to read only
voltage at this frequency (actually, at a double frequency, but this is
immaterial for our general discussion), all the other frequencies being
filtered out. Moreover, it will \emph{look} into a specified \emph{phase} of
this signal. The exact meaning of this last statement will be explained below
and in Appendix~\ref{app:ac}.

The \emph{measuring input} (actually, there are two of them, A and B, and you
can also hook the instrument up to measure the difference $\mathrm{A}-\mathrm{B}$)
is an AC voltmeter, which will measure rms values, and where the sensitivity of
a readout can be changed according to experimental needs, from several volts to
a few millivolts full scale. It can be connected across different parts of a
circuit: for example, it can be connected across $R$ or across $C$. Let us now
discuss specific measuring arrangements: one for measuring a resistance,
another for measuring a capacitance.

\subsection{How to use lock-in to measure resistance?}

Suppose we want to measure a resistance, which may be changing in the course of
a measurement, for example, as a function of temperature. This will be the case
in our \VOtwo\ measurement in this lab. We may connect the lock-in as shown in
Figure~\ref{fig:lockinR}.

\begin{figure}[htbp]
  \centering
  \tikzfig[width=0.8\linewidth]{fig-lockin-resistance}{%
    Circuit diagram. A box labeled Lock-in amplifier sits at the top, with two
    terminals for the measuring input, across which V-meas appears, and two for
    the oscillator voltage output, which supplies V-zero RMS. From the
    oscillator output a wire runs straight down through the resistor R-zero and
    then through the resistor R-D to the right-hand outer contact of the sample.
    The sample is drawn as a horizontal strip carrying four contacts, labeled
    from left to right I-plus, V-plus, V-minus and I-minus. The two inner
    contacts are connected straight up to the measuring-input terminals; a
    parasitic capacitance C-p, drawn with dashed lines, bridges those two inner
    leads. The return wire leaves the second oscillator terminal, runs down the
    right-hand side, along the bottom of the figure and up the left-hand side to
    the left-hand outer contact of the sample, so the current path forms one
    closed loop through the sample.}
  \caption{Connecting the lock-in amplifier for a four-probe resistance
    measurement. $R_0 = \SI{600}{\ohm}$ is the oscillator's internal output
    resistance, $R_D$ the limiting decade resistor, and $C_p$ the parasitic
    capacitance of the cables. The current enters and leaves through the outer
    contacts $I_{\pm}$; the voltage is measured across the inner contacts
    $V_{\pm}$.}
  \label{fig:lockinR}
\end{figure}

Looking at the front panel of the lock-in reveals that the oscillator circuit
has internal resistance $R_0 = \SI{600}{\ohm}$, which is really a part of the
lock-in, but, for clarity, is shown on the diagram above outside of a lock-in
block. It is connected in series with the limiting decade resistance $R_D$ and
the sample. The measuring input of the lock-in is connected across the inner
contacts of the sample.

\paragraph{Let us first discuss the current.}
It is alternating, AC current; we will be always talking about rms values of all
voltages and currents in this discussion. The limiting resistor $R_D$ serves to
limit the current, and to make this current approximately constant. To do so,
$R_D$ should be much greater than the sample resistance $R_{S\,tot}$, which is
the total resistance of a sample between the \emph{outside} current contacts
shown in the figure above, and also includes the resistance of these current
contacts (a contact to a sample may have appreciable resistance, see
Appendix~\ref{app:fourprobe}); it is the total resistance encountered by the
current as it flows through the sample.

The value of $R_{S\,tot}$ will be changing in the course of a measurement, as we
change the sample temperature. Let us say its maximum value is
$R_{S\,tot} = \SI{500}{\kilo\ohm}$, and the minimum value can be neglected
compared to \SI{500}{\kilo\ohm}, which is indeed the case with \VOtwo\ (assuming
contact resistance is not very large). Suppose we decided that we will be
satisfied with keeping current constant to \SI{5}{\percent}, i.e.\ we are not
worried about up to \SI{5}{\percent} possible overall distortion in our
measurement. The maximum and minimum rms currents then will be
$I_{\max} = V_0/(R_0+R_D)$ when $R_{S\,tot}$ is negligible, and
$I_{\min} = V_0/(R_0+R_D+R_{S\,tot})$ when $R_{S\,tot}$ is maximum; here $V_0$
is the rms voltage at the oscillator output. We want to have
\begin{equation}
  \frac{\Delta I}{I_{\min}}
    = \frac{I_{\max}-I_{\min}}{I_{\min}} = 0.05 ,
\end{equation}
or, as it is easy to see, simply $R_{S\,tot}/(R_0+R_D) = 0.05$, or
$R_0 + R_D = R_{S\,tot}/0.05 = \SI{10}{\mega\ohm}$. Compared to this value,
$R_0 = \SI{600}{\ohm}$ can be neglected, and essentially our limiting resistor
$R_D$ should be set at \SI{10}{\mega\ohm} (which, incidentally, is the maximum
resistance which can be set in our decade box). If we use this large limiting
resistor, the current will be equal to
$I = \SI{2}{\volt}/\SI{10}{\mega\ohm} = \SI{0.20}{\micro\ampere}$.

\paragraph{Let us now discuss the voltage signal.}
Small current $I$ will result in a fairly small voltage $I R_S$, especially when
$R_S$ also becomes small. Lock-in can measure small voltages, but below about
the \SI{1}{\milli\volt} level, the precision will suffer. If we want to
increase the current, in the present scheme we will have to sacrifice some
additional accuracy in our overall measurement of $R_S$ over the whole range of
its variation.

The sample has four contacts: two outside contacts for the current and two
inner ones for the voltage. You should read about the four-probe measurement in
Appendix~\ref{app:fourprobe}, and a brief description of this principle should
be given in your Final Report. The voltage from the inner leads
$V_{\mathrm{meas}}$ (see the figure above) is supplied to the lock-in's
measuring input, where it will be amplified (thus the name lock-in
\emph{amplifier}), and provided for the readout on the front panel and, through
the digitizing interface, to the computer. The resistance between the inner contacts will be found as
$R_S = V_{\mathrm{meas}}/I$. Above, when we were discussing the current, we
considered $R_{S\,tot}$ between the outside current contacts, but now we are
measuring resistance $R_S$ between the inner voltage contacts. The lock-in will
display the RMS value as a DC-level signal.

The maximum voltage $V_{\mathrm{meas}}$ in our example (measuring \VOtwo) is
going to be found at room temperature; it will correspond to the maximum value
of $R_S$. Suppose the maximum $R_S$ is equal to \SI{300}{\kilo\ohm} (note that
we are using here $R_S < R_{S\,tot}$). Then the maximum
\begin{equation}
  V_{\mathrm{meas}} = I R_S
    = (\num{2e-7})\times(\num{3e5}) = \SI{60}{\milli\volt}.
\end{equation}
Such a signal would be quite sufficient, but the trouble is that it will be a
lot lower once $R_S$ decreases at high temperatures, where, above the
semiconductor-to-metal transition (see below), we may get a signal as small as
about \SI{0.1}{\milli\volt}. The lock-in is going to give us poor resolution in
that range. Recall also that, with our proposed numbers, we will have a
distortion in resistance values of about \SI{5}{\percent} over the whole range,
as our current will increase by \SI{5}{\percent} when $R_S$ drops by such a
large factor.

An alternative way of dealing with changing $R_S$ is to measure the current
instead of stabilizing it; it can be best measured by monitoring voltage across
the known limiting resistor. Then we do not need such a large limiting
resistor: instead of being \SI{10}{\mega\ohm} it can be, for example,
\SI{10}{\kilo\ohm}, and the current will be 1000 times higher, i.e.\
\SI{0.200}{\milli\ampere} instead of \SI{0.200}{\micro\ampere}. But we will have
to monitor the current and to use the measured current $I_{\mathrm{meas}}$ in a
calculation of a sample resistance at every point along the way. This can be
done with the use of an additional computer-interfaced voltmeter, and with
modification of a LabVIEW program.

\notebox{Do not overdrive the sample}{The current should not be increased
indefinitely; excessive currents can produce self-heating and even damage the
thin-film sample. Do not use currents greater than \SI{1}{\milli\ampere} with
our \VOtwo\ samples.}

\paragraph{So, how is the lock-in measuring resistance?}
The oscillator output, as we said above, provides sinusoidal voltage of rms
value $V_0$ at a fixed stable frequency which we can set at will. The readout is
tuned to read only voltage at this frequency (actually, at a double frequency),
all the other frequencies being filtered out. Additionally, it will read only a
specified phase of this signal. For our resistive measurement we expect that the
specified phase should be zero, as resistor does not produce a phase shift. So,
if we are measuring pure resistance, we should be interested in the
\emph{in-phase} signal (see Appendix~\ref{app:ac}). The circuit and the phasor
diagram of a pure resistor $R_S$ are shown in Figure~\ref{fig:resistorphasor}.

\begin{figure}[htbp]
  \centering
  \tikzfig[width=0.95\linewidth]{fig-resistor-phasor}{%
    On the left, a series circuit: an AC source of RMS voltage V-zero drives a
    current I through three resistors in series, R-zero, R-S and R-D. On the
    right, a phasor diagram with a horizontal in-phase axis and a vertical
    out-of-phase axis. A single heavy arrow lies along the in-phase axis and is
    labeled V-R-S equals I times R-S equals V-zero R-S divided by the sum of
    R-S, R-zero and R-D. A curved arrow to the right of the diagram marks the
    counterclockwise rotation with angular frequency omega.}
  \caption{Circuit and phasor diagram for a purely resistive sample.}
  \label{fig:resistorphasor}
\end{figure}

As always, the phasor should be imagined rotating counterclockwise with
$\omega$; the in-phase direction rotates with it; the projection onto the
horizontal axis is the real signal. The current is determined by the source
voltage $V_0$ applied across all of the resistors. We see that in any other
direction the signal will be smaller than in this in-phase direction. Thus, if
we adjust the phase till we see the \emph{maximum} voltage at the measuring
input, this maximum voltage will correspond to the true value of resistance
according to $R_S = V_{\mathrm{meas}}/I$. Simultaneously, the
\ang{90}-shifted, out-of-phase signal should be zero. It is easier to tune the
out-of-phase signal to zero than to find the true maximum in-phase, because when
measuring the near-zero signal we can increase the sensitivity. So the usual
procedure is to change the phase till one brings out-of-phase signal as close to
zero as possible and then to shift the phase by \ang{90} displaying the in-phase
signal. Further, all of the above should be independent
of frequency $f = \omega/2\pi$. This is what one expects; however, there are
complications.

\paragraph{Complications} arise because, in addition to the resistance, the
circuit contains the so-called \emph{parasitic} or \emph{stray} capacitance
$C_p$. This is the capacitance of shielded cables, etc. In our circuit it is
typically of the order of \SIrange{10}{100}{\pico\farad}. This parasitic
capacitance can be thought of as being plugged in parallel to the sample
resistance $R_S$, as shown in the circuit diagram above where it is drawn in
dashed lines. It will effectively shunt (short out) the resistance, leading to
incorrect, decreased values of $V_{\mathrm{meas}}$.

\paragraph{Solution.}
Recall that lock-in allows you to adjust the oscillator frequency $f$, which of
course implies that reactance $X_C = 1/(2\pi f C_p)$ can be controlled. The
solution is then clear: in order to minimize the detrimental effect of $C_p$ we
need to decrease $f$ as much as possible: at low $f$ the reactance
$1/(2\pi f C_p)$ will be large, and the shunting effect of the parasitic
capacitance will be minimized. On the other hand, using extremely low frequency
$f \sim \SI{1}{\hertz}$ may be difficult, as the lock-in tends to lose
its ability to lock the phase at extremely low frequencies. Additionally, the
so-called flicker or $1/f$ noise increases at low frequencies. But we do not
need to reduce $f$ that much: using $f \approx \SI{50}{\hertz}$ is comfortable
in terms of this phase-locking and achieves what we need. Indeed, at
$f = \SI{50}{\hertz}$, for $C_p = \SI{100}{\pico\farad}$, we find
$X_C = 1/(2\pi f C_p) \approx \SI{30}{\mega\ohm}$, which is more than sufficient
to prevent a decrease in $V_{\mathrm{meas}}$ due to shunting.

\notebox{Practical advice}{Insert a \emph{known} resistor in place of your
sample (preferably, with the value similar to $R_S$) and empirically find the
frequency range at which your measurement agrees with that known resistor
value. This is an example of a \emph{baseline} measurement which is done to
establish the validity of the measurement technique.}

\subsection{How to use the lock-in amplifier to measure the capacitance?}

The sample is now a capacitor; the material we want to study (\BTO) is a
dielectric inside the capacitor. In this measurement we do not need to worry
about additionally limiting the current; the current we are dealing with here
is sufficiently limited by $R_0 = \SI{600}{\ohm}$. Here we will be using lock-in
as a device to measure signals that are \emph{phase shifted} from the phase of
$V_0$ due to the presence of resistors and capacitors in the circuit.
The basic circuit diagram for the capacitance measurement is shown in
Figure~\ref{fig:lockinC}.

\begin{figure}[htbp]
  \centering
  \tikzfig[width=0.8\linewidth]{fig-lockin-capacitance}{%
    Circuit diagram. A box labeled Lock-in amplifier sits at the top, with two
    measuring-input terminals, across which V-meas appears, and two
    oscillator-output terminals supplying V-zero RMS. From one oscillator
    terminal a wire runs down through the resistor R-zero to point 1, the top
    plate of the sample capacitor. The capacitor, marked sample and drawn as two
    plates with a shaded dielectric between them, has terminals numbered 1 and 2
    at the top and 3 and 4 at the bottom; the sample is inserted by connecting
    leads 2 and 4 between points 1 and 3. From point 3 a wire returns down and
    round to the second oscillator terminal. The two measuring-input leads run
    down the left of the figure to points 1 and 3, so the measuring input sits
    across the capacitor, and a parasitic capacitance C-p, drawn with dashed
    lines, is connected in parallel with the sample.}
  \caption{Connecting the lock-in amplifier for a capacitance measurement.
    $R_0 = \SI{600}{\ohm}$ is the built-in internal resistance of the signal
    generator, and $C_p$ is the parasitic capacitance of the empty circuit,
    which acts in parallel with the sample.}
  \label{fig:lockinC}
\end{figure}

As in a resistance measurement described above, the signal is generated by the
lock-in amplifier's internal reference (oscillator) output; in this measurement
we will be setting it at $V_0 = \SI{2.0}{\volt}$ (actually, at
\SI{1.999}{\volt}), $f = \SI{25}{\kilo\hertz}$. As before, the built-in internal
resistance of the signal generator is $R_0 = \SI{600}{\ohm}$, and there exists
parasitic capacitance $C_p$ shown in dashed lines; this capacitance can be
thought of as connected in parallel to sample capacitance.

Surprising as it may sound, we will first do the measurement of an \emph{empty}
circuit, without the sample. Later we will insert the sample by connecting
sample leads~2 and~4 between points~1 and~3, as shown in the diagram above. The
purpose of the measurement without the sample is in tuning the lock-in, so as to
make it later correctly measure the sample's capacitance.

First, leaving the sample unconnected, we will tune the lock-in so that it is
in-phase with the oscillator; it will measure $V_0 = \SI{2.0}{\volt}$. We then
change the phase by \ang{90} and look at $V_{\mathrm{meas}}$ while gradually
increasing the sensitivity of the measuring readout. On a sensitive-enough scale
we will see some non-zero signal due to parasitic capacitance of our ``empty''
circuit. We will adjust the phase watching the signal on increasingly more and
more sensitive scales till this out-of-phase signal will be made as close as
possible to zero. So, in other words, we will follow the same procedure as in
measuring resistance, as was described above. See Appendix~\ref{app:ac} for the details of the
phasor diagrams in this case.

Once we zeroed the out-of-phase signal, we will insert the sample, so that
points~2 and~4 are connected to~1 and~3; once we have done that, the out-of-phase
readout from the sample $V_{\mathrm{meas}}$ will appear at the measuring input
and the magnitude of that out-of-phase signal is going to be proportional to the
capacitance of the sample. Indeed, the capacitance of our sample is usually
small, of the order of \SI{100}{\pico\farad}--\SI{1}{\nano\farad}, and so, with
$R_0 = \SI{600}{\ohm}$ and at our chosen frequency, we have $\omega R_0 C \ll 1$.
In this limit the in-phase component of the signal will be very close to the
voltage provided by the signal generator, $V_0$, and the out-of-phase component
will be close to $V_0\,\omega R_0 C$ (see Appendix~\ref{app:ac} for the details).
By measuring the out-of-phase component $V_{\mathrm{meas}}$, knowing $V_0$,
$\omega$ and $R_0$, we can calculate the capacitance of the sample as
\begin{equation}\label{eq:Csimple}
  C = \frac{\delta}{\omega R_0},
  \qquad\text{where}\qquad
  \delta = \frac{|V_{\mathrm{meas}}|}{V_0}.
\end{equation}

\subsection{Lock-in amplifier settings}

The standard settings of the EG\&G model 5209 lock-in amplifier are:

\begin{center}
\begin{tabular}{@{}ll@{}}
  \toprule
  \textbf{Setting} & \textbf{Value} \\
  \midrule
  Input            & A \\
  Sensitivity      & \SI{3}{\volt} \\
  Filter           & flat, track \\
  Oscillator level & \SI{1.999}{\volt} \\
  Oscillator frequency & \SI{25}{\kilo\hertz} \\
  Reference phase  & tuned (see below) \\
  Reference        & internal \\
  Oscillator out   & BNC cable connected \\
  Time constant    & \SI{300}{\milli\second} \\
  Dynamic reserve  & high resolution \\
  Output display   & signal \\
  \bottomrule
\end{tabular}
\end{center}

Before measuring the sample, we do a practice run by measuring the capacitance
of a piece of BNC cable. See Appendix~\ref{app:ac} for the relevant phasor
diagrams.

\begin{figure}[htbp]
  \centering
  \photo[width=0.85\linewidth]{lockin-front-panel.png}{%
    Photograph of the instrument rack. On top is an EG\&G Princeton Applied
    Research model 5209 lock-in amplifier; its front panel is divided into
    sections labeled Sensitivity, Filters, Tuning, Reference and Output, with a
    digital display reading 6.0, an analog meter, and a second display reading
    0.00. A BNC cable is plugged into the oscillator output at the right. Below
    it is a Keithley 199 System DMM/Scanner whose red display reads 001.090
    kilohms.}
  \caption{The lock-in amplifier (top) and the Keithley 199 multimeter
           (bottom).}
  \label{fig:frontpanel}
\end{figure}

First we check the signal generated by the internal source of the lock-in.
Connect one end of a BNC cable to the ``osc out'' connector, and the other end
to input~A. In the ``tuning'' section scroll until the display shows the
reference phase, ``ref ph'', and set it to \num{0.0}. The analog meter and the
output display on the right should show a voltage slightly less than
\SI{2}{\volt}.

We are going to be interested in the out-of-phase component, and we can get that
by pressing the \textsc{90 deg} button once. Ideally, the out-of-phase signal
should be zero. However, with all the electronics inside the lock-in, keeping the
phase of the output voltage under control is not an easy job and there will be a
small but finite out-of-phase signal.

The engineers designing the lock-in provided a way to correct any phase error:
run an autotune cycle. First press the red ``auto'' button, and then press the
button slightly up and left. Wait a few seconds, and the lock-in will find the
best phase so that the in-phase signal is maximum, and the out-of-phase signal
is zero. To check the out-of-phase, press the \textsc{90 deg} button again and
you should see zero.

Upon closer inspection you may discover that the autotuning is not perfect. Change
the sensitivity from \SI{3}{\volt} to \SI{300}{\milli\volt}, and you will see
that there is a finite signal. You should do a manual fine tuning of the phase,
so that this signal is as close to zero as possible. You may increase the
sensitivity to \SI{100}{\milli\volt} and continue tuning the phase.

\notebox{Do not overload the input}{Do not increase the sensitivity further,
because the input of the lock-in will be overloaded. If you see the red overload
sign flashing, change to a less sensitive setting. Never set the lock-in to
sensitivity of \SI{10}{\milli\volt} or higher.}

Now that you have the out-of-phase voltage at zero, select the sensitivity of
\SI{300}{\milli\volt} again. Insert a BNC ``T'' connector to the line and
connect another BNC cable. Suddenly, there will be a finite out-of-phase signal
$V_{\mathrm{meas}}$. Calculate $C$ using
$C = (V_{\mathrm{meas}}/V_0)/(\omega R_0)$ with $V_0 = \SI{2.00}{\volt}$,
$\omega = \SI{1.57e5}{\per\second}$ and $R_0 = \SI{600}{\ohm}$. Measure the
length of the cable. Test a longer cable.

\notebox{Question}{In round numbers, what is the capacitance of one foot of BNC
cable?}

% =============================================================================
\section{The rest of the experimental apparatus in detail}\label{sec:apparatus}
% =============================================================================

\subsection{Sample chamber}

A photograph of the sample chamber is shown in Figure~\ref{fig:chamber}. There
are four wires reaching the inside of the chamber, ``red'' and ``black'' pairs.
The sample will be either a film of \VOtwo\ with four resistive leads, or a
capacitor with \BTO\ requiring only two leads attached to its plates. If two
students do the measurement, it is recommended that each one of them mounts
their own sample and evaluates the data independently. However, \VOtwo\ samples are
fragile, and a staff member should be asked to mount them in the chamber, at
least initially.

\begin{figure}[htbp]
  \centering
  \photo[width=0.62\linewidth]{sample-chamber.png}{%
    Photograph looking down into the open sample chamber. A cylindrical aluminum
    block, wrapped on the outside with many turns of copper heater wire, has a
    circular well in the middle. Inside the well lie small samples with red and
    black wires soldered to metal clips. A platinum resistor for temperature
    control is visible at the upper right, and the edge of the aluminum cup
    cover appears at the top right of the frame.}
  \caption{The sample chamber, opened. A heater is wrapped around the aluminum
    block; the platinum resistor for temperature control is visible at the upper
    right. A second platinum resistor (barely visible) sits in the center of the
    chamber and is used to measure the temperature of the samples.}
  \label{fig:chamber}
\end{figure}

In normal use, the sample chamber is covered by an aluminum cup (see the upper
right corner of the photograph). The whole assembly is covered by a
double-walled stainless steel cover for thermal insulation. A double-walled
evacuated vessel is called a \emph{dewar}. If you wish to cool the sample faster,
remove the stainless steel dewar. You may also remove the aluminum cup, but in
that case stop taking data, since the temperature readout will not be reliable.

\subsection{Samples and sample preparation}

Samples of \VOtwo\ in a form of thin films (about \SI{100}{\nano\metre} thick)
have been deposited on Si/SiO\textsubscript{2} substrates in another lab on campus.
Contacts to them will be attached by staff members. These samples are rather
fragile, especially those contacts, and should be handled with great care. Ask a
staff member to plug the sample into the holder, at least initially.

The \BTO\ sample is from a commercial distributor. It comes in ceramic form with
many different, randomly oriented microcrystals. Feel free to split off your
sample from larger chunks of the ceramic blocks. Pick a flat and thin piece of
ceramic \BTO, and shape it to a thin plate by lightly sanding on the sides. Keep
in mind that capacitance $C = \varepsilon A/d$: the thinner the better, and the
larger the area the better. Also remember that you will need to estimate $A$ and
$d$ numerically, and so the more regular the sample shape, the easier and more
accurate this estimate will be.

\begin{figure}[htbp]
  \centering
  \photo[width=0.62\linewidth]{silver-paint.png}{%
    Photograph of the sample-preparation bench. On the left stands an open glass
    vial of grey silver paint; behind it is the paint container labeled High
    Purity Silver Paint. In the center of the bench lie a small black two-pin
    connector and a small pale chip of ceramic barium titanate. On the right is
    the vial cap with a long thin wire used to apply the paint.}
  \caption{Materials for mounting the \BTO\ sample: silver paint, a two-pin
           connector, a ceramic chip, and the thin wire used to apply the
           paint.}
  \label{fig:silverpaint}
\end{figure}

Use silver paint to create the capacitor plates and mount the sample (with the
silver paint) onto the two-terminal connector shown in the middle of the figure.
The best way to do this is by placing the sample between the two pins and using
the silver paint to create the capacitor plates. The silver paint should be
applied by a piece of thin wire (visible in the photograph above). Be careful
not to paint over the edges, because that creates a short circuit and ruins the
sample. Remember that you are making a capacitor.

\notebox{The paint dries very quickly}{Do not leave the silver-paint container
open for an extended period of time; open it for as short an interval as
possible, just to get the paint onto a piece of wire.}

\begin{figure}[htbp]
  \centering
  \photo[width=0.48\linewidth]{two-pin-connectors.png}{%
    Photograph of three small black two-pin connectors side by side. The one on
    the right is empty, with two bare gold pins. The one in the center has a
    translucent chip of ceramic barium titanate mounted between its pins. The one
    on the left has a small test capacitor soldered across its pins.}
  \caption{An empty two-pin connector (right), a mounted \BTO\ sample (center),
    and a \SI{100}{\pico\farad} test capacitor soldered to the connector (left).}
  \label{fig:connectors}
\end{figure}

A stereo microscope is provided for mounting the sample, although with good eyes
and eye--hand coordination it can be done without the microscope.

\notebox{Return your sample}{The sample of \BTO\ (but not of \VOtwo) you make
for the studies should be returned with your final lab report. Simply attach the
sample with Scotch tape to the first page of your report.}

\begin{figure}[htbp]
  \centering
  \photo[width=0.6\linewidth]{bnc-connections.png}{%
    Photograph of the bench-top wiring. The bottom of the aluminum sample
    chamber stands at the top of the frame on a white insulating base, with a
    white Teflon-insulated cable and red and black wires running from it down to
    a metal panel carrying four BNC bulkhead connectors. One BNC is fitted with a
    black coaxial cable and another carries a short-circuit termination.}
  \caption{Connecting the apparatus. One BNC cable runs from the signal source to
    input~A of the lock-in, and another from the lock-in to the BNC leading to
    the sample. When the sample is in place, the other BNC must be terminated
    with a short circuit to ground, as shown.}
  \label{fig:bnc}
\end{figure}

Connect the experimental apparatus: one BNC cable from the signal source to
input~A of the lock-in, and another cable from the lock-in to the BNC leading to
the sample. Check that there is no sample in the sample chamber. (For later use,
when you put in the sample, you need to terminate the other BNC with a short
circuit to ground, as shown in the photograph.)

The out-of-phase signal you see now is due to all of the capacitances leading to
the sample, including the capacitance between the Teflon-insulated black wire and
ground. Do a fine tuning of the phase so that this signal is zero. Switch to a
sensitivity of \SI{300}{\milli\volt} and insert the sample. The signal you see
now is entirely due to the sample, and you can use
$C = (V_{\mathrm{meas}}/V_0)/(\omega R_0)$ again to calculate the capacitance of
the sample. (If the signal is small, you may need to change the sensitivity to
\SI{100}{\milli\volt}.)

For all of this to work, $V_{\mathrm{meas}}/V_0 = \omega R_0 C$ should be much
less than~1. In Appendix~\ref{app:ac} we derived the exact formula for the
calculation of the capacitance, and we noted that the difference between the
exact and the approximate formulae is about \SI{1}{\percent} when
$\delta = 0.1$.
Therefore, for $V_0 = \SI{2}{\volt}$, the systematic error caused by using the
approximate formula will be less than \SI{1}{\percent} if $V_{\mathrm{meas}}$ is
smaller than \SI{200}{\milli\volt}. If you made a very thin sample and the
capacitance is large, you may go outside of this limit. When that happens you should
either use the exact formula or reduce the frequency to \SI{10}{\kilo\hertz} or
\SI{5}{\kilo\hertz} and repeat the tuning.

At this point you are ready to measure the temperature dependence of the
capacitance. Keep in mind that we expect the capacitance to increase as we warm
the sample, and select a range where the current signal is less than one third of
the full range. This way you do not have to change sensitivity when the most
interesting part of the measurement is happening.

\subsection{Temperature control}

The ``guts'' of the temperature controller~\cite{ref:pid} are shown in
Figure~\ref{fig:tcguts}. The ``brain'' is an Omega CNi8DH44 unit, connected to a
solid-state relay (on the left) and other circuitry, like fuses and connection
pads. The solid-state relay converts the low-power output of the converter into
the high power that drives the heater wrapped around the sample chamber (see the
photograph of the sample chamber above). The input for the temperature controller
is a platinum resistor that is placed in the aluminum block of the sample
chamber. The resistance of platinum (and all other metals) depends on the
temperature, and this resistor has a resistance of \SI{100}{\ohm} at temperature
\SI{0}{\degreeCelsius}. For
platinum, a \SI{1}{\degreeCelsius} increment of temperature changes the
resistance by \SI{0.38}{\ohm}.

\begin{figure}[htbp]
  \centering
  \photo[width=0.72\linewidth]{temperature-controller.png}{%
    Photograph of the inside of the open temperature-controller enclosure. A
    black Omega CNi8DH44 module with a white specification label sits at the top
    right. To its left are a solid-state relay and a transformer, and the base of
    the box is filled with blue, grey, black and yellow wiring running to a row
    of screw terminals along the front edge.}
  \caption{Inside the temperature controller.}
  \label{fig:tcguts}
\end{figure}

The simplest way to warm the sample to a pre-set temperature and keep it there is
the ``on/off'' control mechanism. You specify the \emph{set point} (or
\emph{setpoint} in one word), which is just the highest temperature needed, and
the controller starts to heat the sample chamber with full power. If the ramp is
specified, the controller will slow down this heating rate by lowering the power
to keep the specified ramp. This is important in order to obtain good data. You
should experiment with the ramp rate and settle at the acceptable rate which you
will judge by the quality of your data.

Once the temperature at the control thermometer reaches the set point, the
heating power is switched off. It is most likely, however, that at this moment
there is an amount of excess heat in the system, and temperature keeps climbing
(this is called an ``overshoot''). After a while the readout of the control
thermometer starts decreasing, and crosses the setpoint from above. The
``on/off'' controller switches ``on'', but by this time the rest of the sample
chamber is too cold, and the temperature keeps dropping for a while. Then the
temperature overshoots again (in the opposite direction) and the process goes
into a cycle. The desired temperature will never be reached; instead there will
be an oscillation around it. To fix this problem the temperature controller is
programmed to use ``PID'' control, with the P, I and D parameters optimized for
this particular sample chamber. Read about the PID control in
Ref.~\cite{ref:pid}.

\notebox{Question}{Be ready to answer this question: what is the role of the
``I'' in PID?}

\notebox{Hot surfaces and overheating}{Most of the capacitive measurement is done
at temperatures over \SI{100}{\degreeCelsius}. Be careful: the sample chamber
will be hot. Do not exceed \SI{160}{\degreeCelsius}, because the sample chamber
will be damaged. The sample chamber can be seriously overheated if left alone
with the heater on. Always switch off the temperature controller when the
measurement is done. If you are using ``ramp'' mode, remember to make sure that
``soak'' is also on for a short time, to ensure that the heater shuts off once
the desired temperature is reached; more on this in
Appendix~\ref{app:tempcontrol}.}

\subsection{Temperature measurement}

In measurements of temperature-dependent properties it is a good practice to
have two thermometers: one to control the temperature (described above) and
another one to measure the actual temperature of the sample. This second thermometer is often
directly attached to the sample. In our case, for the sake of easy sample
changes, the thermometer is simply inserted into the center of the sample
chamber. This thermometer is also a platinum resistor.

We use a Keithley 199 multimeter in four-probe resistivity mode to read out the
resistance of the temperature probe. Learn about the four-probe method in
Appendix~\ref{app:fourprobe}.

\notebox{Question}{Be ready to answer this question: what is the advantage of the
four-terminal method over the two-terminal method, and how is this advantage
achieved?}

The multimeter is set to ``Ohm'' mode and the four leads should be connected to
the four terminals. (Note the two red leads go to the same side of the resistor,
and the two black leads go to the other side. For the extension wires the reds
are connected to browns and the blacks are connected to whites.)

When the computer reads out the resistance values, it will convert them to
temperature according to the (approximate) relation
\begin{equation}
  T = \SI{100}{\degreeCelsius}
      + \frac{R - \SI{100}{\ohm}}{\SI{0.385}{\ohm\per\degreeCelsius}} .
\end{equation}

The two data files contain (1) the time dependence of the temperature and (2) the
temperature dependence of the lock-in output voltage. It is expected that you
transfer the relevant data files to a thumb drive to carry out the evaluation of
the data.

Before measuring the sample, study the temperature control process. Start at room
temperature or at a temperature below \SI{50}{\degreeCelsius}. Select a set point
of \SI{70}{\degreeCelsius}. Record the time dependence of the temperature. Put
the temperature controller into ``on/off'' mode and keep recording. Temperature
oscillations will develop. Put the temperature controller into ``PID'' mode.
Discuss these observations in your lab report.

% =============================================================================
\section{Suggested protocol for a measurement}
% =============================================================================

\begin{enumerate}
  \item Tune the lock-in, and set up the sample (\VOtwo\ or \BTO) in the sample
        chamber. Cover the sample chamber with the stainless steel cover (the
        dewar).
  \item Start the \texttt{Ferro.vi} LabVIEW routine.
  \item Start saving data in LabVIEW.
  \item Program ``Set point 1'' of the temperature controller to
        \SI{90.0}{\degreeCelsius} for a \VOtwo\ measurement and
        \SI{150.0}{\degreeCelsius} for the capacitance \BTO\ measurement.
  \item Program the ramp rate. Change the ramp rate and examine the data: you
        want to make sure that your ramp rate is slow enough for the data to be
        reliable. Discuss with the laboratory staff what this means.
  \item Watch the signal as the sample crosses the phase transition. In \VOtwo\
        this will take place over a fairly wide temperature range; in \BTO, in
        a narrow temperature range.
  \item Once the setpoint is reached you will take more data while the sample is
        cooling.
  \item If you are not in ramp/soak mode, program the lower set point of the
        temperature controller to \SI{10.0}{\degreeCelsius}.
  \item The sample will slowly cool while the data is recorded. In the \BTO\
        measurement, nothing interesting happens below \SI{80}{\degreeCelsius},
        so when the temperature reaches \SI{80}{\degreeCelsius} you can take off
        the dewar to speed up the cooling. When measuring a \VOtwo\ sample, you
        want to take the data in a slow-enough ramp mode in the whole
        temperature range.
\end{enumerate}

% =============================================================================
\section{Advice for evaluating the capacitance data}\label{sec:advice}
% =============================================================================

From the graph of the measured capacitance as a function of temperature you will
find the $T_C$ for this material.

Estimate the area of the capacitor and the distance between the plates and
calculate the relative dielectric constant $k$ at room temperature. To find $k$,
you will need to calculate the capacitance $C_0$ of an empty (i.e.\ filled with
vacuum or air) capacitor with the same geometry as your sample, and then take
your measured $C$ and divide it by that calculated $C_0$: $k = C/C_0$.

Give an estimate of the error in $k$ due to the necessarily inaccurate
measurement of the area and the distance between the painted plates. This error
is expected to be large. Compare the value of $k$ to the known dielectric
constant of this ceramic material.

Our main interest, however, is to use the data from our ceramic sample to access
some of the intrinsic, fundamental properties of \BTO. This means that your
values of $k$ will need an important correction, which we discuss here. If the
dielectric constant is very large, the capacitance of the sample should be very
large too. However, in this measurement, the capacitance will not increase near
the transition that much. A possible explanation is that in a ceramic sample the
dielectric polarization is interrupted at the grain boundaries. The capacitance
between the grains provides coupling, but that inter-grain capacitance does not
change with temperature, and certainly does not diverge at $T_C$.

You can model this behavior by taking the sample as \emph{two capacitors in
series}: $C_b$ (capacitance of the grain boundaries) and $C_{BTO}$ (capacitance
of the bulk grains of \BTO). For the effective (measured) capacitance
in this case we have
\begin{equation}
  C = \left[C_b^{-1} + C_{BTO}^{-1}\right]^{-1}.
\end{equation}
Here $C_b$ is independent of the temperature and $C_{BTO}$ corresponds to an
ideal \BTO\ uniformly-filled capacitor. In this model you can assume that at the phase transition
point, when $C_{BTO}$ is very large, the $C_{BTO}^{-1}$ term can be neglected,
and all you measure is $C_b$, i.e.\ $C \approx C_b$. Thus, from your measurement
of the peak capacitance at $T_C$, you will know $C_b$. For all the other
temperatures you can calculate $C_{BTO}$ as
\begin{equation}\label{eq:CBTO}
  C_{BTO} = \frac{C_b\,C}{C_b - C} = \frac{C}{1 - C/C_b},
\end{equation}
where $C$ is the measured capacitance and $C_b$ is the maximum (peak)
capacitance at $T_C$. Therefore, you should use this corrected $C_{BTO}$ and not
the measured capacitance $C$ directly in finding $k$ values: $k = C_{BTO}/C_0$.

Once you obtain corrected $k$ versus $T$ data, refer to
Appendix~\ref{app:physics} (which you should be reading continually while doing
this experiment; it contains most of the physics), analyze Curie--Weiss
dependence of $k$ in the vicinity of the transition, and find Curie--Weiss
constant $C_{CW}$. From it, calculate the microscopic dipole moment and the
displacement of Ti ion in the unit cell based on your data. Compare it to the
displacement you will find in the literature (you will have to find this
value, which is known from X-ray measurements of \BTO). Also calculate the same
quantities using the intrinsic (single-crystalline) value of $C_{CW}$ from the
literature.

\clearpage
\appendix
% =============================================================================
\section{Elements of AC circuit theory relevant to our measurements using the
         lock-in amplifier}\label{app:ac}
% =============================================================================

We are dealing with sinusoidal AC voltages and currents, such as
$V = V_0\cos(\omega t + \gamma)$, where $V_0$ is the amplitude, $\omega$ is the
angular frequency, and $V(t=0) = V_0\cos\gamma$ is the initial condition, i.e.\
a signal at $t = 0$; $\gamma$ is called the \emph{phase angle}, or the
\emph{phase}. You will recall that on a \emph{phasor diagram} which you studied
in your first physics course, all signal amplitudes are represented as rotating
vectors; $\omega$ is the angular velocity of a rotating phasor related to
regular frequency by $\omega = 2\pi f$; projections of these vectors onto a
chosen axis are real instantaneous signals. A \emph{cosine} function used above
corresponds to a projection onto a horizontal axis; if we used a \emph{sine}
function, it would represent a projection onto a vertical axis. In
Figure~\ref{fig:phasorrot} we show $V_0$ at the moment of time $t$, when it makes
an angle $\omega t + \gamma$ with the horizontal; we see $V_0$'s projection onto
the horizontal direction, $V_0\cos(\omega t + \gamma)$, which is the real signal
at the moment $t$.

\begin{figure}[htbp]
  \centering
  \tikzfig[width=0.85\linewidth]{fig-phasor-rotating}{%
    Phasor diagram. A dashed circle is centered on the origin of a pair of faint
    axes. A heavy arrow labeled V-zero runs from the origin to the circle at an
    angle above the horizontal; the angle between the arrow and the positive
    horizontal axis is marked omega t plus gamma. A dashed vertical line drops
    from the tip of the arrow onto the horizontal axis, and the segment of the
    horizontal axis from the origin to that foot is drawn heavy; a callout
    identifies it as V equals V-zero cosine of omega t plus gamma, the projection
    of V-zero onto the horizontal direction, that is, the real signal at the
    moment t. A curved arrow above the circle indicates counterclockwise rotation
    with omega.}
  \caption{A rotating phasor and its projection onto the horizontal axis.}
  \label{fig:phasorrot}
\end{figure}

We will recognize that a phase angle also specifies a \emph{direction} on a
phasor diagram; there is an infinite number of possible radial directions in a
circle.

It is convenient to measure all phases with respect to the phase of the source,
$V_0$. A zero phase angle then corresponds to the direction of $V_0$; we will
call it in what follows the \emph{in-phase} direction. A direction perpendicular
to the in-phase is called \emph{out-of-phase}. It is clear that projections of
$V_0$ onto the in-phase and out-of-phase directions are $V_0$ and zero
respectively.

\begin{figure}[htbp]
  \centering
  \tikzfig[width=0.85\linewidth]{fig-phasor-directions}{%
    Phasor diagram. A dashed circle is crossed by two perpendicular grey lines
    that are tilted with respect to the page: one is labeled in-phase direction
    number one, running from lower left to upper right, and the other is labeled
    out-of-phase direction number two. A heavy arrow labeled V-zero runs from
    the center along the in-phase direction to the circle. A curved arrow above
    the circle indicates counterclockwise rotation with omega.}
  \caption{The in-phase and out-of-phase reference directions, defined relative
           to the source voltage $V_0$.}
  \label{fig:phasordir}
\end{figure}

Deciding to use these two phases (or these two phase directions), we are
effectively setting our initial $\gamma$ to zero. But we will be eventually
measuring average rms values which do not depend on the initial phase; at the
same time, all \emph{relative} phases will be preserved in this new system of
phase referencing.

Measuring a specified phase of a signal (or ``looking'' at a phase, as we say in
electronics) is equivalent to looking at a signal projection onto a chosen
direction in a phasor diagram. While we chose two such directions, in-phase and
out-of-phase with respect to the source voltage $V_0$, the lock-in is capable of
looking at any phase direction. This is one of its advantages over a normal
voltmeter.

With this, let us now consider a phasor of an RC circuit.
Figure~\ref{fig:rcphasor} shows a phasor of a series connection of $R$ and $C$,
with AC source supplying voltage $V_0$; it is a vector sum of voltages $V_R$ and
$V_C$. For completeness we also show the current $I$, which is established in the
circuit; it is shown pointing in the direction of $V_R$. The current is shown
only once because it is the same for all elements in a series connection.

Note that current is said to be \emph{in-phase with the voltage across the
resistor}, $V_R$. We should be careful not to mix our in-phase direction along
$V_0$ as it was defined above with the in-phase direction for the current; these
are two different phase directions, as can be seen on the diagram.

\begin{figure}[htbp]
  \centering
  \tikzfig[width=0.95\linewidth]{fig-rc-phasor}{%
    On the left, a series RC circuit: an AC source supplying V-zero and current I
    drives a resistor R and a capacitor C connected in series. On the right, a
    phasor diagram whose in-phase and out-of-phase reference axes are drawn
    tilted. Three heavy arrows leave the common origin: V-R points up and to the
    right, V-C points down and to the right at right angles to it, and V-zero
    lies between them, completing the rectangle drawn with dashed construction
    lines. A longer arrow along the direction of V-R is labeled I.}
  \caption{Series $RC$ circuit and its phasor diagram.}
  \label{fig:rcphasor}
\end{figure}

It is clear that projections of $V_R$ and $V_C$ onto both in-phase and
out-of-phase directions are non-zero, while the algebraic sum of projections of
vectors $V_R$ and $V_C$ onto the in-phase direction is $V_0$ and to the
out-of-phase direction is zero. Projections of $V_R$ and $V_C$ constitute signals
measured with the lock-in amplifier across $R$ or $C$ respectively, either
in-phase or out-of-phase.

Measuring signals in a particular phase as described above is called
``phase-sensitive'' detection. We are not discussing here how it is practically
implemented in the inside electronics of a lock-in amplifier, but only its
meaning in terms of a phasor diagram.

Measuring resistance is discussed in the main text in
Section~\ref{sec:techniques}; we need to watch for the detrimental effect of a
parasitic capacitance, and the practical solution is to measure at low-enough
frequency.

\subsection{Measuring capacitance}

Here is an illustration of the lock-in tuning process in terms of phasor
diagrams.

\paragraph{Initial tuning.}
The source voltage $V_0$ is directly measured, $V = V_0$. The phase is adjusted
so that there is no out-of-phase component (Figure~\ref{fig:tune1}).

\begin{figure}[htbp]
  \centering
  \tikzfig[width=0.62\linewidth]{fig-tune-initial}{%
    Phasor diagram with a horizontal in-phase axis and a vertical out-of-phase
    axis. A single heavy arrow labeled V-zero equals V lies exactly along the
    in-phase axis, so the out-of-phase component is zero.}
  \caption{Initial tuning: the measured voltage lies entirely along the in-phase
           direction.}
  \label{fig:tune1}
\end{figure}

\paragraph{The capacitor is added to the circuit.}
The voltage drop on the resistor $V_1 = V_R$ and the voltage drop on the
capacitor $V_2 = V_C$, via vector addition, add up to the source voltage
$V_0$. There is a \ang{90} phase difference between $V_R$ and $V_C$. The phasor
diagram looks as shown in Figure~\ref{fig:tune2}. The phase angle $\alpha$ is
given by
\begin{equation}
  \tan\alpha = \frac{V_R}{V_C} = \omega R C .
\end{equation}
For $\omega = 2\pi f = 2\pi \times \SI{25}{\kilo\hertz}$,
$R_0 = \SI{600}{\ohm}$ and $C = \SI{100}{\pico\farad}$ we get
$\tan\alpha = \num{0.0094}$. For such small angles (in radians!) we can safely
take $\tan\alpha = \alpha = \omega R C$.

For $V_0 = \SI{2.0}{\volt}$ the measured out-of-phase voltage is
$V_2 = V_C\sin\alpha$. For small $\alpha$, $V_C \approx V_0$ and
$V_2 = -V_0\,\omega R_0 C = \SI{-19}{\milli\volt}$.

\begin{figure}[htbp]
  \centering
  \tikzfig[width=0.66\linewidth]{fig-tune-capacitor}{%
    Phasor diagram with horizontal in-phase and vertical out-of-phase axes. A
    heavy arrow labeled V-C equals V points to the right and slightly downward,
    making a small angle alpha below the in-phase axis. Dashed construction lines
    drop from its tip to the two axes, marking a large in-phase component V-one
    on the horizontal axis and a small negative out-of-phase component V-two on
    the vertical axis.}
  \caption{With the capacitor added, the measured voltage acquires a small
           negative out-of-phase component.}
  \label{fig:tune2}
\end{figure}

\paragraph{The parasitic capacitance, retuned.}
Connecting the measurement circuitry is like adding a capacitor to the circuit.
Here $C = C_p$ represents the sum of all parasitic capacitances. Now we tune the
lock-in so that the out-of-phase component is zero again. The angle $\alpha$ is
still small, so $\tan\alpha = \alpha = \omega R_0 C_p$
(Figure~\ref{fig:tune3}).

\begin{figure}[htbp]
  \centering
  \tikzfig[width=0.66\linewidth]{fig-tune-parasitic}{%
    Phasor diagram with horizontal in-phase and vertical out-of-phase axes. After
    retuning, the heavy arrow labeled V-C equals V lies along the in-phase axis.
    A short arrow labeled V-R points straight up along the out-of-phase axis,
    and the arrow labeled V-zero, the vector sum of the two, lies a small angle
    alpha above the in-phase axis.}
  \caption{After retuning, the measured signal is again purely in-phase; the
           source phasor $V_0$ now sits at the small angle $\alpha$.}
  \label{fig:tune3}
\end{figure}

\paragraph{The sample is added.}
Finally we add the sample. The capacitance gets larger: $C' = C_p + C$, where $C$
is the capacitance of the sample. Since the larger capacitance results in smaller
impedance, the current in the circuit increases, and the voltage drop on the
resistor increases. The phasor diagram now looks as in Figure~\ref{fig:tune4},
where $\tan\beta = \omega R_0 C'$.

\begin{figure}[htbp]
  \centering
  \tikzfig[width=0.68\linewidth]{fig-tune-sample}{%
    Phasor diagram with horizontal in-phase and vertical out-of-phase axes. A
    heavy arrow labeled V-C-prime equals V points to the right and slightly
    downward, at a small angle phi below the in-phase axis. A steeper arrow
    labeled V-R points up and slightly to the right, and their vector sum,
    labeled V-zero, lies at the angle beta above the in-phase axis. Dashed
    construction lines complete the parallelogram.}
  \caption{With the sample inserted, the measured out-of-phase signal is
    $V_2 = V_{C'}\sin\varphi$ with $\varphi = \alpha - \beta$.}
  \label{fig:tune4}
\end{figure}

The measured out-of-phase signal is $V_2 = V_{C'}\sin\varphi$, where
$\varphi = \alpha - \beta$. As long as all of these angles are small, we get
$V_2 = V_0\,\omega R_0 C$, where $C$ is the capacitance of the sample. Therefore
\begin{equation}\label{eq:Cdelta}
  C = \frac{\delta}{\omega R_0},
\end{equation}
where $\delta = -V_2/V_0$ ($\delta$ is a positive quantity and, of course, so is
the capacitance). Here $V_2 = V_{\mathrm{meas}}$.

The capacitance of the typical sample is in the \SI{100}{\pico\farad} range. For
\SI{100}{\pico\farad}, at \SI{25}{\kilo\hertz} and $R_0 = \SI{600}{\ohm}$, the
quantity $\omega R_0 C$ is much less than~1. Therefore, in this limit,
\begin{align}
  V_1/V_0 &\sim 1, \\
  V_2/V_0 &\sim -\omega R_0 C, \\
  C       &= -\frac{V_2/V_0}{\omega R_0} = \frac{\delta}{\omega R_0}.
\end{align}

\subsection{Alternative derivation using complex numbers}

The same result can be obtained using complex numbers instead of phasors.
Complex numbers and symbols for treatment of AC circuits were introduced in the
``Primer to Linear Circuits and Transmission Lines'', posted with the Course
Documents on Blackboard. The impedance of the capacitor is $Z = 1/(j\omega C)$.
The current in the RC circuit shown in Figure~\ref{fig:rccomplex} is
$I = V_0/[R + 1/(j\omega C)]$, and the voltage on the capacitor is

\begin{figure}[htbp]
  \centering
  \tikzfig[width=0.42\linewidth]{fig-rc-complex}{%
    A simple series circuit drawn as a rectangle: a resistor R along the top left,
    a capacitor C along the top right, and an AC source symbol supplying V-zero
    along the bottom.}
  \caption{The series $RC$ circuit used in the complex-number derivation.}
  \label{fig:rccomplex}
\end{figure}

\begin{equation}
  V = \frac{V_0\,/\,j\omega C}{R + 1/j\omega C}
    = \frac{V_0}{1 + j\omega RC}
    = \frac{V_0}{1 + (\omega RC)^2}
      - \frac{j\,V_0\,\omega RC}{1 + (\omega RC)^2}.
\end{equation}
The real and imaginary parts are
\begin{equation}
  V_1 = \frac{V_0}{1 + (\omega RC)^2}
  \qquad\text{and}\qquad
  V_2 = -\frac{V_0\,\omega RC}{1 + (\omega RC)^2}.
\end{equation}
With the proper setting of the lock-in amplifier, described in the main text, we
can measure $V_2$. Since we know $V_0$, $\omega$ and $R$, we can calculate the
capacitance of the sample:
\begin{equation}\label{eq:Cexact}
  C = \frac{1}{\omega R}\,\frac{1 - \sqrt{1 - 4\delta^{2}}}{2\delta},
\end{equation}
where $\delta = -V_2/V_0$ (since the measured voltage is negative, $\delta$ is a
positive number).

So, to summarize, when $\omega R_0 C$ is small ($\ll 1$) we can use formula
$C = \delta/(\omega R)$, but for a very large and thin sample, capacitance will
not be small, and one has to use the complete expression \eqref{eq:Cexact}
derived above. A simple calculation shows that the difference between the exact
and the approximate formula is about \SI{1}{\percent} when $\delta = 0.1$ and
about \SI{5}{\percent} when $\delta = 0.2$.

\clearpage
% =============================================================================
\section{The four-probe method: advantages and limitations}\label{app:fourprobe}
% =============================================================================

A simple resistance measurement is performed either by measuring both the current
through the resistor and the voltage across the resistor (an $I$ and $V$
measurement), as in Figure~\ref{fig:a2f1},

\begin{figure}[htbp]
  \centering
  \tikzfig[width=0.6\linewidth]{fig-a2-two-terminal}{%
    Circuit diagram. Current enters from the left through an ammeter, marked A in
    a circle, and passes through the unknown resistor R-x. A voltmeter, marked V
    in a circle, is connected across R-x between two black dots that mark the two
    terminals.}
  \caption{Two-terminal resistance measurement.}
  \label{fig:a2f1}
\end{figure}

or by measuring a current, supplied by a battery, which flows through an
adjustable resistor and the unknown resistor connected in series, with the scale
being set as to read the unknown resistance (an \emph{ohmmeter} measurement,
which we do not show here). Either way, it is a \emph{two-terminal} measurement:
the unknown resistor is connected to the rest of the circuit at two points, or
two terminals, shown in Figure~\ref{fig:a2f1} as black dots. In this measurement
it is implicitly assumed that \emph{contact resistance} and resistance of the
wires is negligible compared to the resistance $R_x$. In many cases, especially
when dealing with commercially-made resistors, this is an excellent
assumption.

There are however many experimental situations in which this measurement is
complicated by the presence of finite, non-negligible \emph{contact resistance},
and sometimes also by the presence of \emph{lead resistance} (resistance of
electrical connections, wires). Consider for example a measurement performed on a
sample shaped as a uniform-width strip (for example, it may be a thin film
sample). When performing a two-terminal measurement, we place two metal contacts
at the ends of a film strip and attach to them current and voltage leads, as
shown in Figure~\ref{fig:a2f2}.

\begin{figure}[htbp]
  \centering
  \tikzfig[width=0.75\linewidth]{fig-a2-strip}{%
    Diagram of a rectangular film strip with one contact, drawn as a black dot,
    at each end. Current enters from the left through an ammeter marked A and
    leaves at the right; short arrows inside the strip show the current flow
    lines, which spread out around each contact. A voltmeter marked V is
    connected between the same two contacts.}
  \caption{Two-terminal measurement on a sample shaped as a strip, with two
           contacts.}
  \label{fig:a2f2}
\end{figure}

There are at least two problems with this arrangement:

\begin{enumerate}
  \item If there is appreciable contact resistance between the film and the
        contacts (black dots), the equivalent electrical circuit describing this
        measurement will look like the one depicted in Figure~\ref{fig:a2f3},
        where the resistors $R_{\mathrm{cont}}$ indicate the presence of said
        contact resistance. If $R_{\mathrm{cont}}$ is not negligibly small
        compared to $R_x$, the voltage drops across them,
        $2 I R_{\mathrm{cont}}$, will be included in the voltmeter measurement.
        Note that here the full sample current $I$ flows through these
        $R_{\mathrm{cont}}$.
  \item Current flow lines in Figure~\ref{fig:a2f3} may be non-uniform around the
        contacts. This makes the length of a sample ill-determined, and thus will
        introduce an error in the calculation of \emph{resistivity}, even if
        contact resistances were negligible.
\end{enumerate}

\begin{figure}[htbp]
  \centering
  \tikzfig[width=0.8\linewidth]{fig-a2-contact-resistance}{%
    Circuit diagram. Current enters from the left through an ammeter marked A and
    passes through three resistors in series: a contact resistance R-cont, the
    sample resistance R-x, and a second contact resistance R-cont. A voltmeter
    marked V is connected across all three, so that both contact resistances lie
    inside the measured section.}
  \caption{Equivalent circuit showing the presence of contact resistance in a
           two-terminal measurement.}
  \label{fig:a2f3}
\end{figure}

We note that the same un-avoidable problems exist in an ohmmeter two-terminal
measurement.

These problems with the two-terminal connection can be overcome by making a
\emph{four-terminal measurement} described below. This measurement is known under
different names such as the \emph{four-probe method}, a \emph{four-terminal
sensing}, a \emph{four-probe geometry}, or as \emph{Kelvin sensing} (the latter
name is rarely used today; it acknowledges that a similar method was first
described by Lord Kelvin a very long time ago). It is universally used in solid
state physics to perform accurate measurements of true material parameters,
primarily of resistivity $\rho$, as well as in other situations when contact
resistance needs to be excluded from the measurement.

The idea is simple: make current and voltage contacts separately. The current is
measured, and so the results will not depend of current contact resistance. The
two inner voltage leads detect \emph{potential difference} between them. The
current practically does not flow into these inner leads, as the voltmeter
resistance is very high. So again, their contact resistances do not enter into
the picture, except for the case when these resistances are becoming very large,
namely, comparable with the input resistance (impedance) of a voltmeter.

\begin{figure}[htbp]
  \centering
  \tikzfig[width=0.85\linewidth]{fig-a2-four-probe}{%
    Diagram of the four-probe geometry. A rectangular, lightly hatched strip
    sample carries four contacts, drawn as large black dots. The current I enters
    from the left and leaves at the right through lead resistances, passing
    through the two outer contacts, which are labeled a and b. Inside the strip,
    five evenly spaced lines run parallel to its length, marked with arrowheads:
    they fan out from contact a, stay uniform and parallel across the middle of
    the sample, and converge again into contact b, so the current density is
    uniform in the region being measured. A dashed arc near each outer contact
    marks where the flow spreads. The two inner contacts are each connected
    straight up to a probe, and a dashed double-headed arrow across the top,
    labeled V, denotes the voltmeter. A dotted double-headed arrow just above the
    strip, between the two voltage leads, marks the distance L. At the right, a
    pair of dotted arrows above and below the strip marks its width w.}
  \caption{The basic four-probe method.}
  \label{fig:a2f4}
\end{figure}

If the current is $I$ and the potential difference between the inner leads is
$V$, then the resistance of the sample which lies between the voltage leads is
simply $R_S = V/I$. Now we can find the resistivity $\rho$ from the formula
\begin{equation}
  \frac{V}{I} = R_S = \frac{\rho L}{w\delta}
  \qquad\text{or}\qquad
  \rho = R_S\,\frac{w\delta}{L},
\end{equation}
where $w$ is the sample width, $\delta$ is its thickness, and $L$ is the distance
between the voltage leads. Here we assume that we know or can measure sample
geometry (uniform width $w$, uniform thickness $\delta$, distance between voltage
leads $L$).

As we said, this method was designed to exclude from the measurement contact and
wire resistances, and it always achieves this goal once we separate the two sets
of contacts. However, the question of \emph{how accurately} one can find the
absolute value of $\rho$ depends somewhat on the placement and shape of the
middle voltage contacts. If we have four contacts which are essentially large
dots placed along the sample (as is shown in the figure above, and is the case in
this lab when we measure \VOtwo), we will not be able to find $\rho$ exactly. The
reason is that our Figure~\ref{fig:a2f4} is somewhat idealized in that current
may not flow along the sample in such a perfect uniform fashion, being somewhat
disrupted by the middle voltage contacts; additionally, we may have a difficulty
measuring $L$ precisely, as it is essentially known to the uncertainty in the
contact dot size. This is illustrated below in Figure~\ref{fig:a2f5}.

\begin{figure}[htbp]
  \centering
  \tikzfig[width=0.8\linewidth]{fig-a2-imperfections}{%
    Diagram of a strip sample carrying four large contact dots, with a current
    lead entering at each end. Two irregular, wavy lines run from the left-hand
    contact to the right-hand one, marked with arrows to show the current
    direction; they pinch together at each of the two inner contacts and bulge
    apart in between, so the current does not flow uniformly along the strip. A
    solid lead drops well below the strip from each of the two inner contacts.
    Above the strip, a dotted double-headed arrow labeled L with a question mark
    spans the gap between the inner contacts; each of those contacts carries a
    pair of dashed extension lines, one at either edge of the dot, showing that
    the separation can only be read to within the size of a contact dot. A dotted
    arrow drops from the middle of that span onto the top of the strip. Below,
    another dotted arrow points up into the strip and is labeled w with a
    question mark, showing that the width is uncertain too.}
  \caption{Imperfections associated with placing contacts on the sample surface.}
  \label{fig:a2f5}
\end{figure}

We will ask you to roughly estimate resistivity of the \VOtwo\ sample, despite
these imperfections. The thickness $\delta$ will be given to you. You may not
have perfectly uniform $w$, but you can estimate average $w$ between voltage
leads; and you can roughly measure $L$ between them as well.

We note that one can further improve the four-probe method in order to accurately
measure $\rho$ by moving inner voltage leads out of the way of the current and
having well defined sample width $w$. We will not address these improvements
here.

\subsection{How much voltage-contact resistance can be tolerated?}

\textbf{Finally, how much of a voltage contact resistance can be tolerated}
before we run into a problem with this measurement? The answer is that
resistances in the voltage contacts will become important if they are comparable
or exceeding the \emph{impedance of our voltmeter} (in our case, the lock-in
amplifier). The input impedance of a voltmeter can be anywhere from
\SI{10}{\mega\ohm} to a gigaohm; one should consult with the instruction manual
of a particular instrument one is using.

To understand why this is so, recall that a real voltmeter \emph{does} need a
tiny current for its operation. The last statement applies to both analog and
digital voltmeters: whether the current is needed to move the \emph{d'Arsonval
galvanometer} in an old analog meter, or to charge a capacitor in a digital one,
it is fundamentally needed nonetheless. (The only voltage-measuring instrument which
does not need a current is a potentiometer, which is not practical in the type of
measurements discussed here.) That is why all voltmeters have finite --- not
infinite --- impedance $R_{\mathrm{imp}}$ (internal resistance): they must admit
some current. So, a voltmeter is essentially a current-sensing instrument
converted to measure a voltage. It fundamentally comprises a current-sensing
element (such as a galvanometer G) in series with a very large resistor
$R_{\mathrm{imp}}$, as depicted in Figure~\ref{fig:a2f6}.

\begin{figure}[htbp]
  \centering
  \tikzfig[width=0.75\linewidth]{fig-a2-voltmeter}{%
    Circuit diagram. A dashed rectangle labeled Voltmeter encloses a large
    resistor R-imp in series with a galvanometer marked G in a circle. This
    branch is connected across the sample resistor R-x, which lies in the main
    horizontal line. The total current I enters from the left; at the first node
    a small current I-G is drawn up into the voltmeter branch, while the
    remainder, I minus I-G, continues through R-x.}
  \caption{Recalling how a voltmeter works.}
  \label{fig:a2f6}
\end{figure}

We can see that $(I - I_G)R_x = I_G R_{\mathrm{imp}}$, and thus
\begin{equation}
  I_G = \frac{(I - I_G)R_x}{R_{\mathrm{imp}}}
      \approx \frac{I R_x}{R_{\mathrm{imp}}}
      = \frac{V_x}{R_{\mathrm{imp}}},
\end{equation}
where we neglected $I_G \ll I$. The instrument, while measuring $I_G$, is
calibrated to read the voltage across the sample $V_x$, which is proportional to
$I_G$.

Now suppose that the voltage contact resistances are comparable to
$R_{\mathrm{imp}}$. The corresponding electrical circuit showing these unknown
contact resistances $R_{\mathrm{cont}}$ is shown in Figure~\ref{fig:a2f7}.

\begin{figure}[htbp]
  \centering
  \tikzfig[width=0.75\linewidth]{fig-a2-voltmeter-contacts}{%
    The same circuit as the previous figure, but now a contact resistance R-cont
    is inserted in each of the two leads that connect the voltmeter branch to the
    sample. The dashed rectangle labeled Voltmeter still encloses only R-imp and
    the galvanometer G, so the two contact resistances lie in series with the
    small current I-G.}
  \caption{A voltmeter measuring a sample with large contact resistances.}
  \label{fig:a2f7}
\end{figure}

Here we have $(I - I_G)R_x = I_G(R_{\mathrm{imp}} + 2R_{\mathrm{cont}})$, or,
again neglecting $I_G \ll I$,
$I R_x \approx I_G(R_{\mathrm{imp}} + 2R_{\mathrm{cont}})$, or
\begin{equation}
  I_G = \frac{V_x}{R_{\mathrm{imp}} + 2R_{\mathrm{cont}}} .
\end{equation}
Since $R_{\mathrm{cont}}$ is not known, the voltmeter will measure incorrect
$I_G$ and thus it will display an incorrect $V_x$.

Note the difference between this circuit and the one depicted in
Figure~\ref{fig:a2f3}: here a \emph{very small} current $I_G$ flows through
contact resistances; there, in a two-terminal measurement, the \emph{full
current} $I$ was flowing through these resistances. This is the main advantage of
a four-terminal over two-terminal measurement.

Generally it is rather difficult to make good contacts to high-resistivity
semiconductors. Contacts to semiconductors and avoidance of contact problems
constitute an active field of study, with many publications devoted just to this
subject. One characteristic of a good contact is that it is \emph{ohmic}, that
is, it has a linear $I$--$V$ curve, as in Figure~\ref{fig:ivohmic}.

\begin{figure}[htbp]
  \centering
  \begin{minipage}{0.45\linewidth}\centering
    \tikzfig[width=0.7\linewidth]{fig-a2-iv-ohmic}{%
      Current--voltage graph with V on the horizontal axis and I on the vertical
      axis. The characteristic is a single straight line of constant positive
      slope passing through the origin.}
  \end{minipage}\hfill
  \begin{minipage}{0.45\linewidth}\centering
    \tikzfig[width=0.7\linewidth]{fig-a2-iv-nonohmic}{%
      Current--voltage graph with V on the horizontal axis and I on the vertical
      axis. The characteristic is an S-shaped curve through the origin: nearly
      flat near zero voltage, then rising steeply at larger positive voltage and
      falling steeply at larger negative voltage, so the slope, and hence the
      resistance, depends on the current.}
  \end{minipage}
  \caption{Left: the linear current--voltage curve of a good, ohmic contact.
    Right: the non-linear curve typical of a high-resistance, ``bad'' contact.}
  \label{fig:ivohmic}
\end{figure}

One often finds that high-resistance ``bad'' contacts are usually also
\emph{non-ohmic}, having a non-linear $I$--$V$ curve, for example like the one on
the right of Figure~\ref{fig:ivohmic}. As one can see, here not only is
$R_{\mathrm{cont}}$ unknown, but it also is not well-defined, depending on the
current $I_G$ that is being passed through it. If in doubt about a contact,
measure its $I$--$V$ characteristic, and stay away from the ones with a
non-linear $I$--$V$.

\notebox{Practical advice}{When preparing to measure the \VOtwo\ sample, first
check its contacts with a simple two-probe ohmmeter. If any of the pairs from the
four contacts turn out to have high resistance (say, of the order of several
Mohms), report this to the TA or to the Instructor, for you will probably not be
able to measure such a sample with the four-probe method.}

\clearpage
% =============================================================================
\section{Programming the temperature controller}\label{app:tempcontrol}
% =============================================================================

The temperature controller can be programmed from the front panel (see the
manual) or by using its web interface. Front panel programming is cumbersome and
not recommended. There are two ways to access the web interface:

\subsection{If you want to change the setpoint only}

Open ``Shortcut to internet'' on the desktop (Internet Explorer). The home page
is served by the temperature controller (the IP address is
\texttt{192.168.1.200}). An ``I Server'' page will appear, with Device Type
Selection. No need to change anything, just click on ``Update''. On the next
page (``Server Home Page'') click on ``Read Devices''. The current temperature
readout of the controller should appear. Click on ``1.'' right before ``Omega''
on the screen. A login screen will appear, enter \texttt{12345678} for
password. You reach the ``Device Setpoints'' screen.

Change ``Setpoint~\#1'' and press update. The temperature controller should show
the new setpoint in the lower left side of the display. As long as you keep the
window open you can change setpoints many times without going through the login
process.

\subsection{If you want to change other settings}

Click on the ``iSeries and MICRO INFINITY Config'' icon on the desktop. The
opening screen will look like Figure~\ref{fig:tc1}.

\begin{figure}[htbp]
  \centering
  \photo[width=0.78\linewidth]{tc-config-open.png}{%
    Screenshot of the iSeries/MICRO-INFINITY Configuration program. The Setpoints
    tab is selected, showing text fields for Setpoint 1 and Setpoint 2, both
    reading 000.0, a drop-down for the display color in the non-alarm state set
    to Green, and a Set Point Deviation checkbox. A column of buttons down the
    right-hand side is labeled Hardware, Read, Write, Help and Quit, and a row
    of tabs along the bottom includes ID code, Setpoints, Input type, Reading,
    Alarm 1, Alarm 2, LB alarm, Out1, Out2, Ramp/Soak, Analog and Comm.}
  \caption{The configuration program as it opens.}
  \label{fig:tc1}
\end{figure}

Click on ``Read'' and on the next screen click on ``File''. The default settings
of our temperature controller are stored in the
\texttt{temperature\_controller.cfg} file (in the ``My Documents'' folder). Click
on the file and the next screen will be like Figure~\ref{fig:tc2}. If the Model\#
and the options come up correctly, click OK.

\begin{figure}[htbp]
  \centering
  \photo[width=0.78\linewidth]{tc-config-file.png}{%
    Screenshot of the configuration program's file-selection dialogue, listing
    the stored configuration file for the temperature controller together with
    the model number and the installed options read back from that file.}
  \caption{Opening the stored configuration file.}
  \label{fig:tc2}
\end{figure}

Next the software needs to establish connection with the temperature controller.
Click on the ``Comm.'' tab and click on the ``Connect'' button under Connection
Wizard. The IP address should be \texttt{192.168.1.200}. Next, the screen shown
in Figure~\ref{fig:tc3} appears. In this window the RS-232 is the default value,
and you need to change it to RS-485. In a few seconds the connection
should be established, and you need to push ``Finish''.

\begin{figure}[htbp]
  \centering
  \photo[width=0.85\linewidth]{tc-comm-wizard.png}{%
    Screenshot of the Connection Wizard dialogue. It shows the network address
    field containing the controller IP address and a selection between the RS-232
    and RS-485 serial standards, with RS-232 selected by default, plus Back, Next
    and Finish buttons.}
  \caption{The connection wizard. Change RS-232 to RS-485.}
  \label{fig:tc3}
\end{figure}

To select a new setpoint, click on the ``Setpoints'' tab and enter the new value
to ``Setpoint 1''.

To activate the ramp mode click the ``Ramp/Soak'' tab on the main window. The
window shown in Figure~\ref{fig:tc4} should appear. Now you can click ``enable''
for the ramp option. The Ramp option allows you to control the rate of heating in
the sample chamber. For example if you set the ramp to 1 hour then it will take
1 hour to get to your setpoint~1, etc. This is important since you
want to keep your sample thermometer consistent with your cavity thermometer.

\notebox{Always enable soak with ramp}{If you decide to use the ramp option you
\textbf{MUST} enable the soak option as well. The soak option will guarantee that
once your temperature is reached and your soak time is complete that the heater
will turn off.}

\begin{figure}[htbp]
  \centering
  \photo[width=0.78\linewidth]{tc-rampsoak.png}{%
    Screenshot of the Ramp/Soak tab of the configuration program, showing enable
    checkboxes and time fields for the ramp and soak phases associated with each
    setpoint.}
  \caption{The Ramp/Soak tab.}
  \label{fig:tc4}
\end{figure}

After that you upload the change to the controller by pressing ``Write'' and
selecting ``Network'' in the screen shown in Figure~\ref{fig:tc5}. During the
upload process hit OK several times when the software stops for confirmation.

\begin{figure}[htbp]
  \centering
  \photo[width=0.78\linewidth]{tc-write-network.png}{%
    Screenshot of the Write dialogue, offering a choice between writing the
    configuration to a File and writing it over the Network to the connected
    controller.}
  \caption{Writing the configuration back to the controller over the network.}
  \label{fig:tc5}
\end{figure}

To access the PID control parameters click on the ``Output 1'' tab. The next
screen shows the ``Control Type''. Here is the place where you can change to
``On/Off Control'' (Figure~\ref{fig:tc6}).

\begin{figure}[htbp]
  \centering
  \photo[width=0.78\linewidth]{tc-output1-pid.png}{%
    Screenshot of the Output 1 tab, showing the Control Type selector with
    entries for PID control and On/Off control, together with the associated
    proportional, integral and derivative parameter fields.}
  \caption{The Output~1 tab, where the control type is selected.}
  \label{fig:tc6}
\end{figure}

\clearpage
% =============================================================================
\section{Data taking with LabVIEW}\label{app:labview}
% =============================================================================

\begin{enumerate}
  \item Click on the ``Ferro'' shortcut on the desktop.
  \item Press the right arrow in the upper left corner to start the program. The
        communication to the voltmeter and the lock-in should start, but the data
        is not recorded yet.
  \item Write the file names where the data should be stored into the fields on
        the top part of the window. Use this format:
        \texttt{c:\textbackslash StudentData\textbackslash 2012Fall\textbackslash
        Group1\textbackslash yourfile.txt}. The directories \texttt{Group1},
        \texttt{Group2} and \texttt{Group3} correspond to the first, second and
        third labs in any given semester.
  \item If you wish to record the data, press ``Save''. Two files can be stored:
        one is the temperature dependence of the capacitance, and the other is
        the time dependence of the temperature. It is crucial that you write the
        file name in your lab notes, where you also record the lock-in settings
        and any other information you may need to reconstruct the circumstances
        of the measurement.
  \item Stop data taking by pressing ``Save'' again. Write in a new file name, so
        that the next time you save data it is recorded separately. (If you do
        not change the file name, the new data will be appended to the existing
        data.)
  \item To clear the screen, press stop. When you restart the measurement (right
        arrow in the upper left corner) a blank screen will appear.
\end{enumerate}

\clearpage
% =============================================================================
\section{Physics of phase transitions}\label{app:physics}
% =============================================================================

\subsection{Introduction to phase transitions}

An interesting, important and useful class of phenomena goes under the name
\emph{phase transition}~\cite{ref:gittrman,ref:papon}. Phase transitions are
widespread in nature, and a number of them find useful applications in
technology. In this lab we will primarily concentrate on phase transitions which
manifest themselves in electrical properties of solids.

Microscopic theories of different phase transitions involve diverse phenomena
specific to these transitions. Yet there are some similarities which allow their
classification and a (somewhat) unified description. The phenomenological theory
--- i.e.\ a non-microscopic theory, one which mathematically describes given
phenomena without asking detailed questions about their mechanism --- of phase
transitions was developed largely by a famous Soviet physicist Lev Landau around
1937~\cite{ref:gittrman,ref:papon,ref:landau1,ref:landau2}. It is based on
thermodynamic ideas and on symmetry. While the detailed description of this
theory exceeds the level at which we can present the subject here, some of the
Landau ideas are actually quite simple and transparent, and we will present them
here, specifically in Section~\ref{sec:landau}.

We will start with a general principle of phase transition classification. Then,
the following sections will detail two phase transitions: the metal--insulator
transition~\cite{ref:dobro,ref:mottbook,ref:mott1949} and the ferroelectric
transition~\cite{ref:jona,ref:blakemore,ref:umich} that you will study in this
lab.

\subsection{Classification of phase transitions}

A phase transition is usually defined as an abrupt change in physical properties
upon a \emph{continuous change in external parameters}. Examples of external
parameters are temperature, pressure, magnetic and electric fields. Of these
temperature is the most common cause of a phase transition, and in this lab you
will be studying such temperature-induced transitions. The value of the
temperature, pressure, or other physical quantity at which a phase transition
occurs is called a \emph{transition point}. If it is temperature, it is usually
denoted as $T_C$.

A distinction is made between two general classes, or, as they are commonly
called, \emph{orders} of phase transitions. This distinction follows pioneering
classification of phase transitions proposed in the 1920s by Paul Ehrenfest. Why do we speak of a transition order? The reason is that Ehrenfest
classified phase transitions based on the behavior of the thermodynamic free
energy~\cite{ref:gittrman,ref:papon} (see below) as a function of other
thermodynamic variables. Under this scheme, phase transitions were labeled by
the lowest derivative of the free energy that is discontinuous at the transition.
Thus, according to this scheme, \emph{first-order} phase transitions are defined
as transitions exhibiting a discontinuity in the first derivative of the free
energy with respect to some thermodynamic variable, and \emph{second-order} phase
transitions are continuous in the first derivative of the free energy with
respect to some thermodynamic variable, but exhibit discontinuity in a
\emph{second derivative}. Free energy itself is always continuous at the
transition point, which is required by conservation of energy.

A modern classification is only slightly different from Ehrenfest's. The meaning
of transition's order in terms of discontinuous free energy derivatives remains,
while the distinction between first and second-order transitions is formulated as
follows: in a first-order phase transition, specified amount of heat, called the
\emph{heat of transformation}, or \emph{latent heat}, is released or absorbed per
unit mass. In a second-order phase transition, there is no latent heat --- heat
is neither released nor absorbed --- but instead some physical quantity, called
an \emph{order parameter} appears; order parameter is equal to zero on one side
of a transition point (in other words, it does not exist in one of the phases),
and it gradually increases from zero on the other side. It
was Lev Landau, in his above-mentioned theory of second-order phase
transitions~\cite{ref:gittrman,ref:papon,ref:landau1,ref:landau2}, who first
raised the order parameter to prominence. We will say more on this subject later.

Examples of first-order phase transitions include many of the most familiar ones,
such as the water--ice transition, or the water--vapor transition; more
generally --- all the solid/liquid/gas transitions are of first order. An
interesting and practically important first-order phase transition manifests
itself in the electrical and optical properties of a solid: it is a
\emph{metal--insulator transition}~\cite{ref:dobro,ref:mottbook,ref:mott1949},
or, as it is sometimes called, a semiconductor-to-metal transition. It can be
found in a fairly large number of solid compounds (but never in pure elements),
including vanadium dioxide (\VOtwo), the material you will be studying in this
lab.

Second-order phase transitions also include many of the transitions of great
interest in physics and its applications. A familiar and important example is a
\emph{ferromagnetic transition}~\cite{ref:gittrman,ref:papon,ref:blakemore,ref:ashcroft}
at a Curie point: above the Curie temperature (in iron it is
$T_C = \SI{770}{\degreeCelsius}$) the material is in a paramagnetic state, with
randomly oriented magnetic moments; at the Curie point the material forms domains
with magnetic moments all aligned in the same direction within a domain. This
phenomenon is called \emph{spontaneous magnetization}. The order parameter we
mentioned above in this case is the net magnetization inside a domain; it is zero
at $T_C$ and increases from zero below $T_C$. We note also that zero value of
order parameter corresponds to a completely disordered state (random orientation
of magnetic moments), while its finite value corresponds to an appearance of
order (common direction for all the magnetic moments in a domain). At the same
time, the material in paramagnetic state is more symmetrical than in the state of
spontaneous magnetization: rotational symmetry of the macroscopic collection of
randomly-oriented microscopic magnetic moments in a paramagnetic state is broken
when material transitions into the oriented state below $T_C$. Thus such a phase
transition is intimately related to symmetry, the fact which was fully explored
by Landau in his theory.

We note here that the ferroelectric transition in barium titanate (\BTO) which
you will study in this lab is similar to the ferromagnetic one, with electric
dipole moments and electrical polarization taking the place of their magnetic
counterparts (we will discuss the differences in Section~\ref{sec:ccw}). We also
note that ideas of spontaneous symmetry breaking migrated from solid state
physics into cosmology and elementary particle theory, and are presently of
singular importance in these fields, illuminating a large number of fundamental
puzzles; for example, the celebrated Higgs mechanism responsible for the
appearance of mass was first proposed by a famous solid state theorist
P.~W.~Anderson in 1962, who at the time worked on the aspects of theory of
superconductivity.

Another example of a second-order phase transition is a \emph{superconducting
transition}~\cite{ref:ashcroft}, which takes place below a transition temperature
(also called the critical temperature) $T_C$: at $T > T_C$ one observes a normal,
resistive metal; at $T = T_C$ the electrical resistance abruptly (in pure metals)
disappears and a normal metal turns into a superconductor with zero resistance;
at lower temperatures $T < T_C$ the resistance is identically zero. The order
parameter in this case is the \emph{superconducting gap}: an energy gap which is
equal to zero at and above $T_C$ and opens up at lower temperatures. It is the gap
in the electronic spectrum, of the order of a meV, separating paired electronic
state (so-called Cooper pairs) from the normal electron spectrum. Note that
superconducting gap is quite different from the gap in a semiconductor or
insulator electronic spectrum, of which we will speak in
Section~\ref{sec:materials}. The entropy of a superconducting state is lower than
that of a normal metal state, indicating that electron pairing provides for a
more ordered state.

\subsection{Thermodynamic approach to phase transitions}

In discussing phase transitions, it is instructive to take a thermodynamic
approach, considering Helmholtz free energy
\begin{equation}\label{eq:helmholtz}
  F = U - TS
\end{equation}
of a material at the two sides of a phase transition. Here $U$ is internal
energy, $T$ is temperature, and $S$ is entropy. Sometimes instead of Helmholtz
free energy $F$, Gibbs free energy $G$ is being considered; it is equal to
$F + pV$, where $p$ is pressure and $V$ is volume. In a phase transition there are
generally two phases, one being stable on one side of the transition, another on
the other side. Thermodynamically, a phase transition occurs at the external
parameter value (temperature, or other) at which \emph{the free energies of the
two phases become equal} (which guarantees, as we said above, that the free energy
is continuous across the transition). One can make some conclusions about the two
general types of phase transitions by making some general observations concerning
formula~\eqref{eq:helmholtz}.

\subsubsection{First-order transitions from the thermodynamic point of view;
               mixed-phase states}

We recall the First Law of Thermodynamics which states that the heat $Q$
absorbed by a system is equal to the change in internal energy of a system
$\Delta U$ plus the work done by a system $W$: $Q = \Delta U + W$. If the system
does not produce any work (i.e.\ if it does not expand or contract against some
outside pressure), then $Q = \Delta U$. This is generally the case in solid-state
transitions of interest to us here. This tells us that first-order phase
transition, which, as we said above, is accompanied by the appearance of latent
heat $Q$, is possible only if internal energy of the system changes. This energy
change must be electronic or ionic in nature: either the electron energy spectrum,
or the electrostatic energy of the ionic lattice, or both, must undergo a change
at the transition point. For example, there may be an energy gap $E_g$
opening up on one side of a phase transition, as is the case in \VOtwo. Such a
rearrangement in the energy spectrum of a solid will be abrupt: in our example,
the semiconductor energy gap in \VOtwo\ exists on one side of the transition and
does not exist on the other; in other words, it does not gradually open up
starting from zero at the transition, as the superconducting gap. The latter, as
we said, serves as an order parameter in that \emph{second-order} transition.
Semiconducting gap, in contrast, could not be considered an order parameter.

Electronic energy spectra in \VOtwo\ on the two sides of a transition are shown
schematically in Figure~\ref{fig:a5f1}.

\begin{figure}[htbp]
  \centering
  \tikzfig[width=0.8\linewidth]{fig-a5-vo2-bands}{%
    Two band diagrams side by side. On the left, labeled semiconductor at T below
    T-C, an empty conduction zone sits above a hatched, filled valence zone, and
    the two are separated by a forbidden energy gap E-g marked with a vertical
    double-headed arrow. On the right, labeled metal at T above T-C, a block of
    empty states sits directly on top of a hatched block of filled states with no
    gap between them. A horizontal dashed line marks the common Fermi level.}
  \caption{A diagram schematically representing the electron energy spectrum in a
    first-order phase transition in which an energy gap opens below $T_C$; such a
    transition is observed in \VOtwo.}
  \label{fig:a5f1}
\end{figure}

Further, we see that in formula~\eqref{eq:helmholtz} there are two terms: $U$ and
$TS$. We already decided that in a first-order transition $U$ must change. But we
also know that the free energies of the two phases must be equal at $T_C$.
Therefore the second term, $TS$, must also change at the transition in order to
compensate for $\Delta U$. Indeed, $S$, the entropy, must change if there is
latent heat $Q$, being related to it by $\Delta S = Q/T$. Change of entropy
signifies appearance or disappearance of order (here the word \emph{order} is used
in a context of order-disorder and not in a context of derivative's order). For
example, in \VOtwo\ one finds different lattice symmetry in the two phases: one
phase is tetragonal, another monoclinic (these terms refer to types of crystalline
atomic arrangements). And indeed one finds a change in entropy which comes
predominantly from this lattice rearrangement (see below a section on \VOtwo). The
lattice constant of a unit cell and material's density are also changing in the
first-order transition. Note that this is different for the second-order
transition, where \emph{order} does not change at $T_C$, but changes gradually
away from $T_C$.

Let us continue with general considerations concerning a first-order transition.
We are considering two phases, phase~1 and phase~2, on the two sides of a
transition. Suppose we are crossing the transition temperature $T_C$ going from
phase~1 to phase~2. The free energy, as we said above, is a continuous function
across the transition, so that $F_1(T_C) = F_2(T_C)$. However, $F(T)$ is not
constant; each phase has lower $F$ on its ``rightful'' side of a transition. This
is schematically shown in Figure~\ref{fig:a5f2}, where phase~1 is stable at
$T < T_C$ and phase~2 is stable at $T > T_C$.

\begin{figure}[htbp]
  \centering
  \tikzfig[width=0.65\linewidth]{fig-a5-free-energy}{%
    Graph of free energy F on the vertical axis against temperature T on the
    horizontal axis. Two smooth curves cross at a single point directly above the
    marked temperature T-C. The curve labeled F-one starts higher on the left,
    falls as temperature rises, and lies below the other curve to the right of the
    crossing; the curve labeled F-two starts lower on the left and rises, lying
    below F-one to the left of the crossing. A dashed vertical line joins the
    crossing point to the label T-C on the temperature axis.}
  \caption{Schematic behavior of the free energy of the two phases across the
           phase transition.}
  \label{fig:a5f2}
\end{figure}

If a given phase --- for example, phase~1, is energetically favored (having lower
$F$), one might expect it to appear right away at crossing the $T_C$. But consider
in more detail the mechanism of its appearance: what has to happen for it to
appear? Some atomic rearrangement and/or lattice symmetry change will be required;
electronic and ionic energy spectra will have to shift; latent heat will have to
be released or absorbed. Very importantly, unless the whole volume of the material
uniformly transitions all at once, domains of phase~1 may form, surrounded by the
phase~2, and, as a result, there appears the \emph{boundary} between the two
phases. This boundary may have energy, either positive or negative. In case
boundary energy is positive, the boundary is said to ``cost'' energy. Appearance
of a boundary thus is changing somewhat the energy balance dictated by the lower
free energy of phase~1. There may appear a \emph{local
energy minimum}, which depends on the local configuration of the phase domains;
this local minimum may preclude the system from falling into the \emph{global
minimum} dictated by the lower $F_1$ value. Thus the phase transition may not
happen exactly at $T_C$, and it may develop in a complex, non-uniform fashion.
During this process, the temperature of the system will stay constant as heat is
added: the system is in a \emph{mixed-phase regime} in which some parts of the
system have completed the transition and others have not.

Such phenomena are well known in first-order phase transitions between solids and
liquids and liquids and vapor: a liquid may be heated to a temperature above the
boiling point or \emph{supercooled} to below the freezing point.
Something similar happens in a first-order phase transition in a solid.
First-order transitions are often characterized by the formation of
\emph{metastable} states. Instead of taking place exactly at $T_C$, the transition
going up in temperature takes place at $T_C + T^{*}$, and going down in
temperature at $T_C - T^{*}$. Here $T^{*}$ is called the \emph{coercive
temperature}. This leads to the phenomenon of \emph{hysteresis}. If we fix a
temperature $T$ somewhere in the interval $T_C - T^{*} < T < T_C + T^{*}$, a
metastable state, often in the form of domains, may survive for long (sometimes
indefinitely long) periods of time on the ``wrong'' side of the transition. In
Section~\ref{sec:mixedstate} you will find a qualitative discussion of the reason
for the appearance of a coercive temperature and hysteresis.

In this lab you will study one such first-order phase transition, namely the
semiconductor-to-metal transition in vanadium dioxide (\VOtwo), which also
exhibits a mixed-phase regime. Additionally, the semiconductor phase of \VOtwo\
has high resistivity, while its metallic phase has low resistivity. Thus, as we
explore the mixed-state with electrical resistance measurement, we are
encountering another interesting type of a transition here, called
\emph{percolation}~\cite{ref:percolation} transition. We will briefly discuss
these phenomena when talking about \VOtwo\ in Section~\ref{sec:materials}. In this
lab, when performing measurements on \VOtwo, you will be able to see how, with
increasing temperature, resistivity gradually changes from high to low,
dramatically changing by several orders of magnitude in the process. The material
goes through the metal--insulator transition, a two-phase mixture occurs, and, as
a result, it also goes through the percolation transition.

\subsubsection{Second-order transitions from the thermodynamic point of view;
               elements of Landau theory}\label{sec:landau}

In a second-order transition there is no latent heat, that is, $Q = 0$ and thus,
from the First Law of Thermodynamics, $\Delta U = 0$. Thus an energy spectrum is
not expected to modify abruptly at the transition point in a second-order
transition. At the same time, the general condition $F_1 = F_2$ at $T = T_C$ must
be satisfied. Looking at formula~\eqref{eq:helmholtz} for $F$, this tells us
immediately that entropy (order) cannot change at $T_C$, but has to change
gradually as the temperature deviates from $T_C$. In other words, the order
parameter must be zero at $T_C$, and will increase away from $T_C$, as we said
above. This, among other considerations, allowed Landau to construct his theory of
the second-order phase transitions. Namely, the fact that order parameter is small
in the vicinity of a transition point allowed him to expand free energy $F$ as a
Taylor series in order parameter, which we will call $m$, and then, making general
physical arguments, to extract significant physical results from this
expansion.

\begin{figure}[htbp]
  \centering
  \photo[width=0.26\linewidth]{landau.png}{%
    Black-and-white portrait photograph of the theoretical physicist Lev Landau,
    facing slightly to one side.}
  \caption{L.~D.~Landau (1908--1968), Nobel Prize of 1962.}
  \label{fig:landau}
\end{figure}

Let us follow his lines of thought, at least in part, here. The Taylor expansion
of a continuous function $F(T,m)$ in powers of a small parameter $m$ around point
$m = 0$, written to the 4-th order of $m$, is
\begin{equation}
  F(T,m) = F_0(T,0) + \frac{F'(T,0)}{1!}m + \frac{F''(T,0)}{2!}m^{2}
         + \frac{F'''(T,0)}{3!}m^{3} + \frac{F''''(T,0)}{4!}m^{4} + \dots
\end{equation}
Here primes indicate derivatives with respect to $m$; for example
$F'' = \mathrm{d}^2F/\mathrm{d}m^2$. So far this is just a mathematical formula
based on continuity of $F$ and smallness of $m$. Landau then argued that $F$
should be invariant with respect to a change in sign in $m$: changing $m$ to
$-m$ should keep the energy unchanged. Indeed, in the absence of external field,
if $m$ is magnetization $M$ or polarization $P$, it seems obvious. This is a
symmetry argument we mentioned before. This implies that our Taylor expansion must
not contain odd powers of $m$. Omitting odd terms with $m$ and $m^3$, and calling
the even-term coefficients
$F''(T,0)/2! = \alpha(T)$ and $F''''(T,0)/4! = \tfrac{1}{2}\beta(T)$ (the factor
$1/2$ introduced here following a convention found in most of the literature), we
obtain
\begin{equation}\label{eq:landau2}
  F(T,m) = F_0(T) + \alpha(T)\,m^{2} + \tfrac{1}{2}\beta(T)\,m^{4} + \dots
\end{equation}

We want to relate behavior of $F(T,m)$ to the observed physical phenomena in the
vicinity of a phase transition. \emph{A minimum of $F$ with respect to an order
parameter $m$ physically implies stable (observable) configuration of our system}:
the value of magnetization or polarization observed in a real system will
correspond to the value of $m$ at which $F$ has a minimum. We will find that these
minima are different for $T > T_C$ and for $T < T_C$, on the two sides of a phase
transition. After finding the minima with respect to $m$, we may be able to see
how $F$ depends on $m$ in the vicinity of these minima. Once we decide on the form
of $F(T,m)$ for fixed $T$, we will be in the position to discuss $T$-dependences.
Further, by taking derivatives of $F(T,m)$, one can construct other quantities,
such as entropy, specific heat, etc. Such is our ``program'' of how to explore
physical consequences of expansion~\eqref{eq:landau2}.

To find minima of $F$ with respect to the order parameter $m$, we require that
$\mathrm{d}F/\mathrm{d}m = 0$ and $\mathrm{d}^2F/\mathrm{d}m^2 > 0$: the usual
conditions we always use when finding a minimum of a function. The first one,
$\mathrm{d}F/\mathrm{d}m = 0$, guarantees that we have zero slope, i.e.\ that we
are at the extrema of a function; the second condition,
$\mathrm{d}^2F/\mathrm{d}m^2 > 0$, guarantees that this extremum is actually a
minimum and not a maximum or the inflection point.

Differentiating~\eqref{eq:landau2} we find
\begin{equation}
  \frac{\mathrm{d}F}{\mathrm{d}m}
    = 2\alpha(T)m + 2\beta(T)m^{3}
    = 2m\left[\alpha(T) + \beta(T)m^{2}\right] = 0,
\end{equation}
which gives us three possible solutions for $m$:
\begin{equation}\label{eq:minima1}
  m_1 = 0
  \qquad\text{and}\qquad
  m_{2,3} = \pm\left[-\frac{\alpha(T)}{\beta(T)}\right]^{1/2}.
\end{equation}

The second condition for a minimum is
$\mathrm{d}^2F/\mathrm{d}m^2 = 2\alpha(T) + 6\beta(T)m^2 > 0$. The simplest way to
satisfy this last inequality is to require that at $T > T_C$,
$\alpha(T) = a\,(T - T_C)$ with coefficient $a > 0$, and that $\beta(T)$ is a
positive constant $b > 0$. We note that if we adopt this form for $\alpha(T)$,
then $\alpha(T_C) = 0$.

So, now the expansion looks like this:
\begin{equation}\label{eq:landau4}
  F = F_0(T) + a\,(T - T_C)\,m^{2} + \tfrac{1}{2}b\,m^{4} + \dots
\end{equation}
And $F$ has minima at
\begin{equation}\label{eq:minima2}
  m_1 = 0
  \qquad\text{and}\qquad
  m_{2,3} = \pm\left[\frac{a}{b}\,(T_C - T)\right]^{1/2}.
\end{equation}

Clearly the expression for $m_{2,3}$, by virtue of containing a square root, works
only for $T < T_C$ (temperatures below the transition); thus $m_1 = 0$ corresponds
to $T > T_C$ (temperatures above the transition). As can be seen
from~\eqref{eq:minima2}, at the transition temperature $m = 0$ in both cases.

Now we will try to explore the physical consequences. Recalling the physical
meaning of $m$, we see that equilibrium magnetization (or polarization) is zero at
$T > T_C$: this is clearly the description of a paramagnetic (paraelectric) state,
in which magnetization from randomly-oriented magnetic moments equals to zero. According to the expansion, at $m = 0$, $F = F_0(T)$. At
$m \ne 0$, and with positive $T - T_C$ (above the transition),
expansion~\eqref{eq:landau4} predicts that $F$ will rise from the lowest energy
level $F_0$ as $F - F_0 = Am^2 + Bm^4$, where $A$ and $B$ are positive
coefficients. This is shown schematically in Figure~\ref{fig:a5f3}(a), where we
took $F_0$ to be zero.

\begin{figure}[htbp]
  \centering
  \tikzfig[width=0.95\linewidth]{fig-a5-landau-wells}{%
    Three graphs of free energy F against order parameter m. In panel (a), for
    fixed T above T-C, the curve is a single smooth well with one minimum at m
    equals zero. In panel (b), for fixed T below T-C, the curve is a symmetric
    double well with two minima at m-two and m-three, placed symmetrically either
    side of a local maximum at m equals zero. In panel (c), for a fixed T below
    T-C but closer to T-C than in panel (b), the same double-well shape appears
    but the two minima are much shallower and lie much closer to m equals zero,
    so the curve is almost flat over a wide region around the origin.}
  \caption{Schematic representation of the dependence of the free energy $F$ on
    the order parameter $m$ at three fixed temperatures.}
  \label{fig:a5f3}
\end{figure}

The minima corresponding to $m_{2,3}$ describe the behavior of $F$ below the
phase transition: our theory predicts that there will be two such minima,
symmetrical in $m$ with respect to $m = 0$. Between these two minima, at $m = 0$,
in virtue of~\eqref{eq:landau4}, we have $F = F_0$ (or zero, since we have chosen
zero at $F_0$). This is depicted in Figure~\ref{fig:a5f3}(b). Physically this
means that our system will have two possible stable states with non-zero
magnetization or polarization. Thus in a ferromagnetic (or ferroelectric) state
magnetization (polarization) will have a \emph{finite} value. This finite
value may be positive or negative; it can reverse, still being consistent with our
theory.

The values of $F$ at these minima can be calculated by substituting $m_1$ and
$m_{2,3}$ into~\eqref{eq:landau4}. They will depend on temperature $T$, which is
now less than $T_C$. We see from~\eqref{eq:landau4} that the closer is $T$ to
$T_C$, the smaller will be the negative term $a(T - T_C)m^2$ in the expansion, and
thus the shallower will be the two minima below the transition. At the same time,
with decreasing $T - T_C$, the two minima will move in towards the $m = 0$ value.
In the limit $T \to T_C$ the two minima will disappear, transforming the
double-minima curve into a single minimum curve (a). At $T = T_C$, because of the
$m^4$ term, the curve will have a significant flat region in the vicinity of
$m = 0$. Physically this must imply that, near the transition point, magnetization
(polarization) in a ferromagnetic (ferroelectric) state can be very easily
reversed. In other words, the material at a temperature just below the transition
point must be very easily polarizable. This suggests that \emph{dielectric
constant} $\varepsilon$ will be high. We will see that indeed this interesting
conclusion is experimentally verified, and you will be able to see this very
clearly in your measurements of \BTO.

So far our conclusions based on the analysis of expansion~\eqref{eq:landau4} have
been very satisfying: we were able to explain some of the most prominent features
of a second-order phase transition. We could continue and explore the physical
consequences for the derivatives of $F$, which will give us entropy and specific
heat. We will find~\cite{ref:gittrman,ref:papon} that entropy $S$, while being
continuous at the transition, decreases below the transition, where there is more
order. We will also find that specific heat experiences a jump at the transition.
Such a discontinuity in specific heat at the second-order phase transition is well
documented experimentally.

One can add external field (magnetic field $H$, or electric field $E$) to this
analysis. In doing so one
finds~\cite{ref:gittrman,ref:papon} that susceptibility $\chi$, which is defined
as $\chi_M = \mathrm{d}M/\mathrm{d}H$ for the magnetic case and
$\chi_E = \mathrm{d}P/\mathrm{d}E$ for the electric
case~\cite{ref:gittrman,ref:papon,ref:ashcroft}, obeys the Curie--Weiss law (here
written for the ferroelectric case):
\begin{equation}\label{eq:landauchi}
  \chi_E = \frac{1}{2a\,(T_C - T)} \ \text{ for } T < T_C,
  \qquad
  \chi_E = \frac{1}{4a\,(T - T_C)} \ \text{ for } T > T_C.
\end{equation}

This shows that Landau theory predicts that susceptibility becomes infinite at
$T = T_C$. We also see that it approaches $T = T_C$ in an asymmetric way
(constants being different by a factor of~2). In an experiment you will perform in
this lab you will see that this is indeed so, except that theoretical infinity is
substituted by a sharp peak.

We recall Curie--Weiss law~\cite{ref:gittrman,ref:papon,ref:ashcroft} which states
that $\chi = C_{CW}/(T - T_C)$ for the paramagnetic (paraelectric) phase, i.e.\
for $T > T_C$. Here $C_{CW}$ is a Curie--Weiss constant (often called Curie
constant). We see that Landau theory identifies coefficient $1/(4a)$ with
$C_{CW}$.

\subsection{Description of the phase transitions and materials you will
            study}\label{sec:materials}

\subsubsection{Metals, semiconductors and insulators, and their conductivity}

With respect to an ability to transmit electrical charge, materials can be
divided into conductors and insulators; this was recognized from the eighteenth
century. Ohm's law, resistance and resistivity were defined and clarified in the
nineteenth and the beginning of the twentieth century. Today, with respect to
electrical resistivity $\rho$ (or its inverse $1/\rho = \sigma$, electrical
conductivity) we distinguish not two but three classes of materials: metals (good
conductors), dielectrics or insulators (very bad conductors), and semiconductors,
which are in some sense intermediate between the first two classes. We also
distinguish them by different ways in which conductivity depends on temperature
(below). True understanding of what is at the core of these different classes came
only in the 1930s, when physicists understood electronic spectra (band structure)
of these materials~\cite{ref:blakemore,ref:ashcroft}. This required creation of
quantum mechanics, formulation of Pauli's exclusion principle, and replacement of
classical statistics with Fermi--Dirac statistics. In Figure~\ref{fig:a5f4} we
show a schematic representation of the electronic bands of a metal, a
semiconductor and an insulator.

\begin{figure}[htbp]
  \centering
  \tikzfig[width=0.9\linewidth]{fig-a5-band-classes}{%
    Three band diagrams side by side against a common vertical axis of electron
    energy, with a horizontal dashed line marking the Fermi level. For the metal,
    on the left, the upper conduction band and the lower valence band overlap in
    energy, and the overlap region is shaded and labeled. For the semiconductor,
    in the middle, the two bands are separated by a small gap straddling the Fermi
    level. For the insulator, on the right, the conduction band and the valence
    band are separated by a much larger gap, marked with a brace labeled
    bandgap.}
  \caption{Electronic spectra (band structure) of metals, semiconductors and
           insulators.}
  \label{fig:a5f4}
\end{figure}

The main distinction between them, which influences electrical conductivity and
its temperature dependence, lies in the presence (in semiconductors and
insulators) or absence (in metals) of an energy gap in the electronic spectrum.
The value of the energy gap is called $E_g$. The highest energy band which
contains states which are completely filled with electrons is called a
\emph{valence band}; an empty band, separated from it by an energy gap, is called
a \emph{conduction band}. Electronic states of the type that may conduct current
(the so-called extended or Bloch states) do not exist in the energy interval
inside the gap.

The difference between semiconductors and insulators is quantitative rather than
qualitative: they both have a gap in their spectra, but the gap in a semiconductor
is smaller than that in an insulator. The rough borderline between these two
classes lies around $E_g = \SIrange{2}{3}{\electronvolt}$. One can say that
semiconductor materials are small band-gap insulators. It should be said that, in
addition to the presence of a relatively small gap, another defining property of a
semiconductor material is that it can be doped with impurities that alter its
electronic properties in a controllable way. Here we will be concerned mainly with
the so-called \emph{intrinsic} semiconductors, which do not have impurity levels
inside the gap. The most famous semiconductors today are Si and GaAs (and their
extended families); but there are hundreds of known semiconductors, and many of
them find useful applications. \VOtwo\ you will study in this lab is one of them.

Let us understand why conduction will be so very different depending on the
presence or absence of a gap in the electronic energy spectrum. You will need to understand
this, as we will ask you to find the value of the gap in \VOtwo\ from your data.

In a conductor, application of an electric field $E$ leads to acceleration of
charge carriers (such as electrons) with mass $m$ and charge $e$: $a = eE/m$. This
acceleration acts for some average time between collisions $\tau$, leading to the
appearance of a drift velocity $v_d = a\tau = (eE/m)\tau$. The current density is
then
\begin{equation}\label{eq:ohm}
  j = e n v_d = \frac{e^{2}n\tau}{m}E = \sigma E
  \qquad\text{(Ohm's law).}
\end{equation}
Here $n$ is the number of carriers per unit volume (carrier density), and $\sigma$
is electrical conductivity which characterizes the material. Its inverse is
resistivity $\rho$.

It is clear from these familiar considerations (constituting part of what is known
as Drude model~\cite{ref:blakemore,ref:ashcroft} of electrical conduction) that an
electron must be able to \emph{change} its energy in order to participate in
electrical current; indeed, it acquires extra kinetic energy $m v_d^2/2$ from the
field; this extra energy will eventually be lost in inelastic collisions of an
electron with phonons (lattice vibrations), other electrons, etc. We know that
eventually this extra energy is deposited into the conducting material, raising
its temperature: the familiar $I^2R$ power dissipated in a conductor comes from
the field (the battery) and ends up as heat in the conducting material. Note that
energy increase $m v_d^2/2$ is, generally, quite small compared to the total
kinetic energy of an electron at the highest filled energy level.

So, to repeat: an electron needs to change its energy in order to accelerate under
the influence of an electric field; this is the prerequisite for electrical
conduction. But, to change its energy, it has to find un-occupied energy levels at
that higher energy. This is possible in a metal, where un-occupied states lie
directly above the occupied ones (at zero temperature we call the highest energy of
occupied states \emph{Fermi level}; it is shown as a dashed line in
Figure~\ref{fig:a5f4}; there are empty states right above it). In a semiconductor
or insulator, in view of the existence of a forbidden energy interval (energy gap)
this is impossible: there are no available states above the top edge of the valence
band (see Figure~\ref{fig:a5f4}). On further reflection, however, we see that this
last statement is strictly true only at $T = 0$. Indeed, at a finite temperature
$T$ electrons may be excited across the gap. Relative probabilities of such
processes can be found from Fermi--Dirac distribution functions; approximate (but
very good) estimates can be done, however, using classical Boltzmann factor
$\exp(-E_g/kT)$. Such an excitation process creates an electron near the bottom of
the conduction band above the gap, and a \emph{hole} (which is equivalent to a
positively charged electron~\cite{ref:papon}) near the top of the valence band.
These charge carriers (electron and hole) are called \emph{mobile carriers}; they
can now participate in conducting the current, because they can find un-occupied
energy levels very close to their respective energies.

The number of mobile electrons and holes depends on the values of $T$ and $E_g$,
and on how many electrons were available for excitation. In good insulators, i.e.\
materials with $E_g > \SI{2}{\electronvolt}$, at practical temperatures (limited,
generally, by the melting point; often we are interested in conductivity at
$T = \SI{300}{\kelvin}$), their numbers are negligibly small. In other words, the
resistivity of an insulator is many orders of magnitude higher than in a
semiconductor or in a metal. For example, resistivity of a very good insulator
amber (the one that ancient Greeks were rubbing to produce static electricity) at
\SI{300}{\kelvin} is \SI{5e14}{\ohm\metre}, while resistivity of pure Si is
\SI{2300}{\ohm\metre}, a difference of eleven orders of magnitude! Metals are
still orders of magnitude more conductive than semiconductors: resistivity of
copper at \SI{300}{\kelvin} is only \SI{1.7e-8}{\ohm\metre}. Thus metals and good
insulators differ by staggering 22 orders of magnitude (!) in resistivity (or
conductivity).

In addition to very different values of $\rho$, its temperature dependencies are
very different as well: in good metals resistivity increases with $T$; in copper
and many other metals this increase is linear at sufficiently high $T$. In
contrast, resistivity in a semiconductor depends on temperature as
\begin{equation}\label{eq:gap}
  \rho = \rho_0\,\exp\!\left(\frac{E_g}{2kT}\right).
\end{equation}
Here the factor of~2 comes from the fact that activation energy is counted from
the Fermi level, or, as it is called in semiconductors, from chemical potential,
which, in intrinsic semiconductors, lies in the middle of the
gap~\cite{ref:blakemore,ref:ashcroft}; an electron has to receive energy equal to
$E_g/2$ rather than $E_g$ to cross the gap. Thus the study of the temperature
dependence of electrical resistance allows us to find $E_g$. The most convenient
way of doing this is to plot $\log(\rho)$ vs.\ $1/T$.

\subsubsection{Metal--insulator transition}

One type of a phase transition is that between metal and insulator or metal and
semiconductor; it is most commonly called \emph{Metal--Insulator
Transition}~\cite{ref:dobro,ref:mottbook,ref:mott1949} or MIT, even when the
poorly-conducting phase is a semiconductor.

If we take a piece of metal, such as copper, and heat it up to its melting point,
the melt will still be a metal. If we boil it (in vacuum) and evaporate the molten
copper, the copper atoms constituting the vapor will be of course, strictly
speaking, forming some sort of an insulator in a sense that they will not be
conducting a current if placed in the electric field. This is not however what we
call a metal-insulator transition. We are interested in a solid material which
retains its general structure and composition, but, under the influence of an
external parameter, for example, temperature undergoes an abrupt transition from a
poorly-conducting semiconductor or insulator state to a well-conducting, metallic
state or vice versa. There are some 50 or more compounds (but no pure elements)
which exhibit MIT. For example, MIT was found in
Fe\textsubscript{3}O\textsubscript{4}, NiS, NbO\textsubscript{2},
V\textsubscript{2}O\textsubscript{3}, etc. One such compound you will have a
chance to study in this lab: it is vanadium dioxide, or \VOtwo. All MIT's are
first-order transitions, although this fact can be obscured by the presence of a
disorder~\cite{ref:dobro}.

\subsubsection{Essential physics of one type of MIT: the Mott transition}

The theory of MITs is rather complex. What is more, several different MIT
mechanisms may be at work in different compounds. Here we can only give a ``taste''
of a relevant theory by qualitatively describing the transition as envisioned by
N.~F.~Mott~\cite{ref:mottbook,ref:mott1949}. Other MIT models
exist~\cite{ref:dobro}.

\begin{figure}[htbp]
  \centering
  \photo[width=0.22\linewidth]{mott.png}{%
    Black-and-white portrait photograph of the physicist Sir Nevill Francis Mott
    in later life, wearing glasses and a suit.}
  \caption{Sir Nevill Francis Mott (1905--1996), Nobel Prize of 1977.}
  \label{fig:mott}
\end{figure}

Sir Nevill Mott, a famous British physicist, in 1949
discussed~\cite{ref:mottbook,ref:mott1949} the following model: consider $N$
hydrogen atoms are in a periodic lattice. If the distances between these atoms are
small, $N$ electrons which they contribute will fill half of the band which
consists of $2N$ states (the reason there will be $2N$ states in such a band is
that there are two states for each electron due to electron spin
one-half: one up and one down~\cite{ref:blakemore,ref:ashcroft}). Thus such a
hydrogen crystal will be a metal. Now Mott imagined increasing the distance between
atoms (increasing lattice constant $a$). With increasing $a$, the band's energy
span will decrease; it will become more and more narrow (consult solid state
physics texts~\cite{ref:blakemore,ref:ashcroft} to see why), but, within the
confines of the band theory, such ``stretched out'' hydrogen crystal will still be
a metal with half-filled band for any, arbitrarily large, separation between
atoms. This is of
course nonsensical, because we know that eventually, at large enough atomic
separations, the crystal must become an insulator.

This leads one to believe that, with increasing $a$, something new has to happen at
a certain distance between atoms. Mott assumed that, starting from some critical
value of $a$, electronic states no longer can be described as extended Bloch
states. Electronic wavefunctions will become localized. Such a crystal will be an
insulator at $T = 0$: each electron will be localized at its respective atom. Such
an insulator is called \emph{Mott insulator}.

At sufficiently high $T$, however, such a material can conduct current, turning
from an insulator into a metal. Indeed, consider two neighboring atomic sites,
1 and 2; each has an electron on it. If these two electrons have equally-oriented
spins, the conduction is impossible in view of Pauli principle: an electron~1
cannot jump to site~2 which is occupied by electron~2 with the same quantum
numbers. But suppose these electrons have opposite spins (what is called
antiferromagnetic order). An electron from an atomic site~1 can in principle jump
to the neighboring site~2, which has an empty energy state for electron~1's
direction of spin (the other state being occupied by the electron with opposite
spin which is sitting on site~2). If electron~1 jumps, it will leave behind a
positively-charged hole on site~1. If such jumps are frequent, with an application
of an electric field the current will flow, this current being carried by both
electrons and above-mentioned holes they leave behind (holes moving in the
direction of the field, and electrons against this direction, as
usual~\cite{ref:ashcroft}). To make this jump, however, an electron has to overcome
Coulomb repulsion which exists between the two electrons residing on the same site
--- it has to have energy at least of the order of this Coulomb repulsion energy.
In a Mott insulator, this Coulomb repulsion barrier energy serves as an analog of a
forbidden gap $E_g$ in a band picture of a dielectric. With raising temperature
more electrons will be excited across this Coulomb gap, and thus the Fermi level
(kinetic energy) will grow. Once the kinetic energy of electrons will become
comparable to Coulomb repulsion energy, the current can flow. Thus, at a certain
sufficiently-high temperature, Mott insulator will undergo a transition to a
metallic state. Similar transitions can be envisioned as a function of doping,
which likewise can increase electron concentration and thus kinetic Fermi energy,
or pressure, which can reduce lattice constant $a$ and, as a result, broaden the
band and raise Fermi energy again.

Therefore, MIT in this picture takes place essentially as a function of
free-electron concentration $n$, at a certain critical value $n_C$. External
parameters which can drive $n$ to this value of $n_C$ can be temperature, pressure,
doping, etc. Lattice transformations can also play a role too, leading to changes
in $n$.

The nature of a MIT in our material, \VOtwo, has been debated for many years. The
latest opinion shared by the majority of workers in the field is that it is
essentially a Mott transition as described above, and not a result of a lattice
transformation. Instead, it is believed that lattice transformation itself is the
secondary effect from an electronic Mott transition.

\subsubsection{MIT in \VOtwo}

The MIT in \VOtwo\ was discovered in 1959, and extensively studied ever since. As
a function of temperature it takes place at
$T_C = \SI{340}{\kelvin} = \SI{67}{\degreeCelsius}$ (transition is sharp in single
crystals~\cite{ref:bergland} but broad in polycrystalline
samples~\cite{ref:gurvitch}; you will measure it in thin films of \VOtwo\ which
exhibit a fairly broad transition). It is a first-order transition characterized
by latent heat, density change, and lattice symmetry change, as well as by a
dramatic change in resistivity. We have no way of detecting aspects of a
transition other than resistivity in this lab.

You will measure resistance as a function of temperature. In
Figure~\ref{fig:a5f5} we show an example of such a measurement taken from
Ref.~\cite{ref:gurvitch}, performed on a \SI{95}{\nano\metre} film of \VOtwo. Your
measurement will be different, as different films exhibit a somewhat different
behavior; for example, the ratio of resistance just below the transition to
resistance above the transition may be quite different in different samples. The
overall features, however, are similar.

\begin{figure}[htbp]
  \centering
  \photo[width=0.8\linewidth]{fig05-vo2-hysteresis.png}{%
    Graph of sheet resistance in ohms, on a logarithmic vertical axis from ten
    cubed to ten to the sixth, against temperature in degrees Celsius from 20 to
    100 on the horizontal axis. Five nearly overlapping hysteresis loops are
    plotted, one for each temperature ramp rate of 20, 10, 5, 1 and 0.1 degrees
    Celsius per minute. Each loop falls from about one megohm at room temperature
    to about 600 ohms above 85 degrees. The cooling branch, marked with a downward
    arrow, drops steeply near 75 degrees; the heating branch, marked with an
    upward arrow, drops steeply near 60 degrees, so the loop encloses roughly 15
    degrees of hysteresis and spans about two and a half decades in resistance. An
    inset plots the absolute temperature shift, from zero to about 0.65 degrees,
    against ramp rate from 0 to 20 degrees per minute at a fixed temperature of 75
    degrees Celsius; the shift rises steadily with ramp rate.}
  \caption{An example of a hysteresis loop in $R(T)$ measured on a
    \SI{95}{\nano\metre} film of \VOtwo\ deposited by pulsed laser deposition
    (PLD), from Ref.~\cite{ref:gurvitch}. The figure shows several loops taken at
    different rates of temperature increase and decrease (temperature ramp rates);
    the upper inset gives a symbol key for the different ramps. The lower inset
    shows temperature shifts $\Delta T$ for different ramps. These are systematic
    errors in measuring temperature. Although on the logarithmic scale of the
    picture all these loops essentially overlap, there is a temperature shift of up
    to \SI{0.6}{\degreeCelsius} for the fastest ramp compared with the slowest.
    \textbf{Note that you will have to worry about such ramp rates and temperature
    shifts in your experiment.}}
  \label{fig:a5f5}
\end{figure}

The transition, as is typical of all first-order transitions (see
Section~\ref{sec:landau}), involves formation of a mixed-state, in which there are
domains of metallic phase and domains of semiconducting phase. This is discussed
in greater detail in Section~\ref{sec:mixedstate}. The transition develops
differently depending on the temperature increasing or decreasing --- there is
substantial hysteresis.

As the temperature increases from room temperature (\SI{25}{\degreeCelsius}),
resistivity will first decrease because our material is a semiconductor. But at
some temperature, which is going to be close to the
$T_C = \SI{67}{\degreeCelsius}$, metallic-phase domains will appear in an
otherwise semiconducting material, and resistivity will start to decrease
faster. Examination of
Figure~\ref{fig:a5f5}, which is plotted on a logarithmic scale in resistance,
suggests that the first deviation from semiconducting behavior indeed takes place
at about \SI{67}{\degreeCelsius}. The fraction of metallic domains grows with
increasing $T$, and at some point a global metallic cluster connects the voltage
leads on a sample; as a result all remaining semiconducting domains will be
shorted out, and resistivity will drop to its metallic value. This is the result
of the percolation transition~\cite{ref:percolation}, which, of course, is made
possible by the appearance and growth of those metallic domains.

\notebox{Question}{If --- for simplicity --- we would imagine that resistivity of a
semiconducting phase is infinitely greater than that of a metallic phase, where, at
what temperature on the curve in Figure~\ref{fig:a5f5}, would be the point at which
a global metallic cluster is forming? This would be the point of a percolation
transition. Read at least some of Ref.~\cite{ref:percolation} and briefly discuss
percolation physics in your report. At what fraction of metal phase is percolation
transition expected to take place in a 2D (thin film) system?}

If the temperature will be now lowered from about \SI{90}{\degreeCelsius} back to
\SI{25}{\degreeCelsius}, the process will reverse, but with substantial hysteresis.
All of this is shown in Figure~\ref{fig:a5f5}. The transition, as can be seen,
spans about 2.5 orders of magnitude in resistance. This measurement was performed
on a \VOtwo\ film; in single crystals resistance can change by 5 (!) orders of
magnitude, and transition is quite sharp, with hysteresis of only about 1--2
degrees~\cite{ref:bergland}.

\subsubsection{Development of a mixed state and hysteresis in
               \VOtwo}\label{sec:mixedstate}

\begin{quote}\small
The text below is taken from Ref.~\cite{ref:gurvitch}, which has been slightly
edited here. In this text you will find a qualitative explanation for the
appearance of the coercive temperature $T^{*}$ and of hysteresis, which was
advanced by the authors of Ref.~\cite{ref:gurvitch}, primarily by Professor Serge
Luryi (Chair, ECE Department, Stony Brook). Other explanations of the same
phenomena were given by different authors, and they can be found in the
literature~\cite{ref:lanskaya}.

The major issue in dispute here is whether each domain exhibits hysteresis on its
own, or whether hysteresis is a collective property of the whole system, arising,
as argued here, as a result of the boundary energy ``cost''. Phases S and M refer
to semiconducting and metallic phases. At the end there is a discussion of a
phenomenon called \emph{non-hysteretic branches}, which you can observe if you do
appropriate measurements on your sample in this lab. Talk to the professor in
charge of this experiment if you want to explore this phenomenon.
\end{quote}

At a given temperature $T$ inside the major hysteretic loop, some parts of the film
have $T_C(x,y) < T$ and others $T_C(x,y) > T$. Here $x$, $y$ are coordinates of a
given location in a film. In the first approximation, the boundary wall between
the S and M phases, at any given temperature $T$, is determined by the condition
$T_C(x,y) = T$. In this approximation, the wall is highly irregular and its
ruggedness corresponds to the scale at which one can define the local $T_C(x,y)$,
i.e.\ to the characteristic length scale of the nanoscopic phase domains. On
closer inspection, however, we need a refinement that takes into account the
\emph{boundary energy}, associated with the phase domain wall itself. The boundary
energy is positive and to minimize its contribution to the free energy the domain
walls are relatively smooth.

Let us examine the process of boundary motion. For concreteness, let us consider
the heating branch. Below the percolation transition, M phase resembles lakes in
the S phase mainland. With raising temperature the area of the M phase increases,
lakes grow in size. When a boundary of a given lake is far from the other lakes,
infinitesimal $\pm\delta T$ changes boundary length by infinitesimal amount, and
the lake area by $\pm\delta A_M$. In other words, when
the lakes are sufficiently separated, we envision a continuous, reversible,
hysteresis-free process of $\mathrm{M}\leftrightarrow\mathrm{S}$ area
redistribution, with neighboring configurations differing microscopically.

Let us now look at the formation of a link between two neighboring regions, which
is the elementary step in the topological evolution of a global percolation
picture. Let us focus on two metallic lakes that are about to merge. Since the
boundary is smooth, at some temperature the distance between the lakes becomes
smaller than the radius of curvature of either lake at the point they will
eventually touch. Therefore, at some $T$ the following two configurations will have
equal energies: one comprising two disconnected M phase lakes that are near
touching, but not quite, and the other with a finite link formed
(Figures~\ref{fig:a5f6}(a) and~\ref{fig:a5f6}(b) respectively).

\begin{figure}[htbp]
  \centering
  \tikzfig[width=0.72\linewidth]{fig-a5-lakes}{%
    Four schematic panels showing a boundary between the semiconducting and
    metallic phases, with the metallic phase shaded. In panel (a) two shaded
    regions approach each other from left and right, their smooth curved
    boundaries almost touching but still separated. In panel (b) the shaded region
    instead runs across the picture as an upper and a lower band joined by a
    narrow neck, that is, a link has formed. Panels (c) and (d), drawn for the
    higher temperature T-zero equals T plus T-star, repeat the same pair of
    configurations but with the lakes grown: in (c) the two boundaries now touch at
    a single point, and in (d) the link between the upper and lower shaded bands
    has become wide.}
  \caption{Semiconductor--metal boundary; the metallic phase is shaded. The top
    row, (a) and (b), corresponds to temperature $T$, and the bottom row, (c) and
    (d), to a higher temperature $T_0 = T + T^{*}$.}
  \label{fig:a5f6}
\end{figure}

Both configurations are characterized by equal boundary lengths and therefore have
equal free energy. The actual transition forming a local link, however, does not
occur at that temperature because of an immense \emph{kinetic} barrier
(meaning: it takes significant atomic rearrangement) between these two
\emph{macroscopically} different configurations. The transition occurs at a higher
$T_0 = T + T^{*}$, when it is actually forced, i.e.\ when the two phases touch at a
point. Here $T^{*}$ is the coercive temperature. In this picture the coercive
temperature arises as a result of having a boundary between different phases; it
does not pre-exist intrinsically within each domain. We associate the steep slopes
of the major loop with the quasi-continuous formation of such links, i.e.\ with
local topological changes, specifically with the merger of metallic lakes on the
heating branch (HB) and of semiconductor lakes on the cooling branch (CB). Near the
major loop ends $T_S$ and $T_M$ the global map consists of widely-separated M and S
lakes, respectively. In these regions we expect to see non-hysteretic behavior, as
was indeed observed.

\paragraph{Non-hysteretic branches (NHB).}
Consider now a small \emph{backward} excursion from $T_0$ on the HB. As the
temperature decreases, some of the M-phase recedes and the S-phase grows, changing
the geometry of the global two-phase map. However, topologically, the last formed
M-link does not disappear immediately for the same kinetic reason. One has two S
regions that need to touch in order to wipe out the M-link. It takes a backward
excursion of amplitude $\Delta T^{*}$ to establish an S-link and thus disconnect
the last M-link. So long as we are within $\Delta T^{*}$, i.e.\ stay on the same
NHB, the area of S and M domains changes continuously, but the topology is stable
and no new links are formed. Within the range of that stable or frozen topology,
$\Delta T^{*} = \Delta T_{\mathrm{NHB}}$, the resistivity of NHB will be
single-valued and its $T$-dependence will be controlled by the percolating
semiconductor phase. This explains the appearance of non-hysteretic
branches and why they have semiconducting slopes.

\subsection{\BTO\ and ferroelectricity}

Ferroelectric materials~\cite{ref:gittrman,ref:papon,ref:jona,ref:blakemore,ref:umich}
exhibit spontaneous electric polarization, the separation of the center of positive
and negative electric charges. Recall the Landau theory conclusion discussed at the
end of Section~\ref{sec:landau}, that the free energy will have a double minimum
below $T_C$: in a ferroelectric state the electric polarization, which is the order
parameter, is non-zero. This creates an excess positive charge on one side of the
material, and negative charge on the other side. The polarization can be reversed
by the application of an appropriate electric field: order parameter can be
flipped from one minimum to the other with E-field. As we said above,
ferroelectricity is named in analogy with ferromagnetism, which occurs in such
materials as iron. The prefix \emph{ferro}, meaning iron, was used to describe the
property despite the fact that most ferroelectric materials do not contain iron.

In most ferroelectric materials the ferroelectricity disappears at high enough
temperature, and the material turns into a regular dielectric. This state is called
the \emph{paraelectric}, in analogy to paramagnets. Most of the known transitions from
the ferroelectric to the paraelectric state are second-order; the transition
temperature is called the Curie temperature, in full analogy with the ferromagnetic
transition. (Note: some ferroelectric materials melt before reaching this
temperature.) In the paraelectric state these materials have very large dielectric
constant: the temperature dependent dielectric constant $\varepsilon(T)$ increases
as the temperature approaches the Curie temperature $T_C$, and shows a divergent
behavior at $T = T_C$. We also note that near $T_C$ the dielectric constant is
essentially equal to the susceptibility, a quantity you will find in many references
on ferroelectricity. The relation between these two quantities is given in the next
section.

Again, recall our --- or rather, Landau's theory's --- prediction made in
Section~\ref{sec:landau}: we predicted very large polarizability near $T_C$, on both
sides of a transition. Large polarizability means that induced charges in a
dielectric will appear easily, compensating free charges on a capacitor plate ---
in other words, that relative dielectric constant $k$ will be high.

Ferroelectricity is often accompanied by significant \emph{piezoelectricity}
(developing electric polarization under mechanical stress) and
\emph{pyroelectricity} (developing electric polarization when heated). Referring to
our favorite (by now) Landau theory, we see that this is perfectly understandable:
the system wants to fall into an energy minimum, and just about any outside stimulus
will induce it to do so.

Ferroelectric materials have many applications, such as multilayer capacitors, gate
dielectrics and waveguide modulators. Pyroelectricity is used in infrared detectors
and other sensors. Piezoelectricity is widely applied in microphones, speakers and
cigarette lighters, and in scientific applications such as the scanning tunneling
microscope.

In this lab you will be studying a well-known ferroelectric material, an oxide of Ba
and Ti, \BTO~\cite{ref:jona,ref:blakemore,ref:doitpoms}. We note that oxides of this
type are important and well studied in solid state physics and materials science,
and not only because of their ferroelectric properties. These oxides all have a
chemical formula of the type ABO\textsubscript{3}, where A is a metal with valency 1
or 2, either A\textsuperscript{+} or A\textsuperscript{2+}; B is a metal with
valency 4 or 5, either B\textsuperscript{4+} or B\textsuperscript{5+}; and
O\textsuperscript{2-} is oxygen. The crystal structure of these materials is called
the \emph{perovskite} structure, after the naturally occurring mineral
CaTiO\textsubscript{3}, which itself is called perovskite. Other examples of
perovskite ferroelectrics include KNbO\textsubscript{3}, PbTiO\textsubscript{3},
KTaO\textsubscript{3} and
PbZr\textsubscript{x}Ti\textsubscript{1-x}O\textsubscript{3} (the latter also called
PZT).

When talking about crystal structure, we often show one chemical-formula unit of
this structure, called the \emph{unit cell}. The whole macroscopic (\emph{bulk}, as
it is often called) crystal of a material can be built by stacking unit cells in all
three spatial dimensions.

\begin{figure}[htbp]
  \centering
  \photo[width=0.85\linewidth]{fig07-perovskite.png}{%
    Two ball-and-stick drawings of a cubic unit cell. In the left-hand cell,
    labeled Perovskites, eight large A atoms sit at the corners of the cube, small
    O atoms sit at the centers of the faces, and a single large dark B atom sits
    exactly at the body center. In the right-hand cell, labeled barium titanate,
    the same arrangement is shown with barium at the corners, oxygen on the faces
    and titanium at the center, but the central titanium atom and the surrounding
    oxygen atoms are displaced slightly along the vertical axis in opposite
    directions, so the cell is no longer symmetric.}
  \caption{(a) Unit cell of a perovskite structure. (b) Ferroelectric distortion in
           \BTO.}
  \label{fig:a5f7}
\end{figure}

Below a certain temperature $T_C$, some of the perovskites are characterized by the
appearance of a \emph{lattice distortion} which changes their structure from cubic,
as in Figure~\ref{fig:a5f7}(a), to a less symmetrical one (the structure which
appears below $T_C$ in our material is called \emph{tetragonal}), as in
Figure~\ref{fig:a5f7}(b). Although all of the atoms constituting the unit cell move
at least somewhat in this distortion, a rough picture of what happens can be
obtained if you will visualize the central Ti\textsuperscript{4+} ion shifting
within the unit cell by a small distance shown to be downwards in
Figure~\ref{fig:a5f7}(b);
simultaneously, the four O\textsuperscript{2-} ions surrounding
Ti\textsuperscript{4+} also shift, but in the opposite direction, i.e.\ upwards in
Figure~\ref{fig:a5f7}(b). All these shifts are quite small in comparison to the
unit cell side $a$; they are somewhat exaggerated in the picture above.

Note that the symmetry of the unit cell is reduced below the transition; a certain
part of the total symmetry which existed above the transition is \emph{broken} below
the transition, as was discussed in Section~\ref{sec:landau}.

When the Ti ion is in the center of the cube, the charges are balanced, and the
polarization is zero. If the Ti and oxygen ions move as illustrated in
Figure~\ref{fig:a5f7}(b) and in Figure~\ref{fig:a5f8}, a finite \emph{electric
dipole moment} $p$ appears in each of the unit cells, and the material becomes
polarized; this polarization can be reversed with an application of an external
field, as shown in Figure~\ref{fig:a5f8}.

\begin{figure}[htbp]
  \centering
  \photo[width=0.9\linewidth]{fig08-batio3-cells.png}{%
    Two three-dimensional ball-and-stick unit cells of barium titanate drawn side by
    side, with a key at the right identifying large red spheres as barium, medium
    purple spheres as oxygen and a small green sphere as titanium. Beside the
    left-hand cell a heavy black arrow labeled P-minus points downward, and the
    green titanium atom sits below the center of the cell. Beside the right-hand
    cell a heavy black arrow labeled P-plus points upward, and the titanium atom
    sits above the center.}
  \caption{Unit cells of \BTO\ with the (green) Ti ion in two off-center positions:
    with $p$ pointing up (called $p^{+}$ here) and with $p$ pointing down (called
    $p^{-}$ here). Below $T_C$, the dipole moments $p^{+}$ and $p^{-}$ appear
    spontaneously without an external field. If there is an electric field $E$, for
    example pointing up, then $p^{+}$ will have lower energy than $p^{-}$.}
  \label{fig:a5f8}
\end{figure}

In this lab we will use a temperature-controlled environment to study the dielectric
constant as the sample undergoes the phase transition, and we will search for signs
of ``divergence'' in the dielectric constant. We will also attempt to extract some
microscopic information from these measurements, as described in the main text of
this write-up and also below.

\notebox{A note on terminology and notation}{The dielectric properties of any
material depend on the frequency at which the properties are measured. This is
described by the dielectric function $\varepsilon(\omega)$. In our case the
properties depend on the temperature as well, and therefore we have a two-variable
function $\varepsilon(\omega,T)$. We use the term ``temperature-dependent dielectric
constant'' to represent the low-frequency limit of the dielectric function,
$\varepsilon(T) = \varepsilon(\omega = 0, T)$. The frequencies applied in our studies
(\SI{25}{\kilo\hertz}) are actually in this low frequency limit. Furthermore, we
will distinguish between dielectric function $\varepsilon(T)$ and dielectric
constant $k(T)$, which is just the relative dielectric function
$k = \varepsilon(T)/\varepsilon_0$. We may call it a ``constant'', following the
normal usage, but in our case, when studying a ferroelectric, we will recognize
that it is a temperature-dependent quantity.}

\subsection{Paraelectric to ferroelectric transition; Curie--Weiss law}

\BTO\ undergoes a second-order paraelectric to ferroelectric phase transition upon
cooling through $T_C = \SI{120}{\degreeCelsius}$. Above that temperature it has
cubic structure and no local dipole moment; below it transforms to a new,
slightly-non-cubic, tetragonal structure in which Ti ion and oxygen ions move in
opposite directions creating a local dipole moment in each unit cell. In other
words, each unit cell becomes electrically polarized and, below $T_C$, these local
dipoles align to form a macroscopic polarization $P$.

As was stated in Section~\ref{sec:landau}, electric susceptibility is defined as
$\chi_E = \mathrm{d}P/\mathrm{d}E$, or, when $P$ is linear with $E$, as
$\chi_E = P/E$. We mentioned in Section~\ref{sec:landau} that Landau theory predicts
that in the vicinity of a phase transition $\chi_E$ is large (in fact, diverging at
$T = T_C$) and that it follows Curie--Weiss Law
\begin{equation}\label{eq:cw}
  \chi_E = \frac{C_{CW}}{T - T_C}.
\end{equation}
The relationship between $\chi_E$ and the dielectric constant $\varepsilon$ is given
by $\varepsilon = \chi_E + 1$~\cite{ref:lockin,ref:accircuits,ref:papon,ref:blakemore}.
Near the transition $\chi_E$ is large, and so, to a very good approximation,
$\varepsilon = \chi_E$. Therefore Curie--Weiss law~\eqref{eq:cw} can be re-written
for the relative dielectric constant as
\begin{equation}\label{eq:cwk}
  k = \frac{C_{CW}}{\varepsilon_0\,(T - T_C)}.
\end{equation}

Again, to clarify the meaning of symbols and to avoid confusion, in this equation
$k$ is a \emph{relative} dielectric constant, in other words, it is equal to the
ratio of capacitance with a dielectric to the capacitance with the same geometry but
with vacuum between the plates, $k = C/C_0$. As was mentioned above, although it is
called a ``constant'', it is not a constant with respect to temperature. The formula
for a parallel-plate capacitance with vacuum (or air) between the plates is
$C_0 = \varepsilon_0 A/d$, where $A$ is one plate's area, and $d$ is the distance
between the plates. The formula for the capacitance with dielectric is
$C = \varepsilon A/d = \varepsilon_0 k A/d$. Here $\varepsilon_0$, the permittivity
of free space, is equal to
\SI{8.85e-12}{\coulomb\squared\per\newton\per\metre} $=$
\SI{8.85e-12}{\farad\per\metre}.

In Section~\ref{sec:landau} we saw that Landau theory predicts this law both above
and below the transition: we had, from~\eqref{eq:landauchi},
$\chi_E = C_{CW1}/(T_C - T)$ for $T < T_C$ and $\chi_E = C_{CW2}/(T - T_C)$ for
$T > T_C$, with $C_{CW1} = 1/(2a)$ and $C_{CW2} = 1/(4a)$, $a$ being expansion
coefficient. This law is indeed observed in the vicinity of ferromagnetic and
ferroelectric transitions, where experimental data mostly concentrates on
paramagnetic or paraelectric states at $T > T_C$. You will have a chance to see it in
your own measurements, and even to test it on both sides of the transition.

\paragraph{Problems.}
As was already discussed in the main text of this write-up, your sample may have a
somewhat irregular shape, so you will at best be able to estimate $k$ perhaps within
a factor of 2. Additionally, you are measuring a ceramic sample, which will tend to
underestimate intrinsic $k$ in your material, as was discussed in
Section~\ref{sec:advice}. The latter problem can be avoided or corrected, at least in
part, by making a correction as described in that section (let us call this a
\emph{capacitance correction}). You want to use your measured $C$ values to calculate
the value of $C_{BTO}$, which is the intrinsic capacitance of your material. Thus
when we will be talking about your capacitance data, we will assume that you made the
suggested capacitance correction. The $k$ values which you will be using in the $k$
vs.\ $T$ data should be using this correction.

Keeping in mind these caveats, by fitting your corrected $k$ vs.\ $T$ data to
formula~\eqref{eq:cwk} you will be able to find (or at least to estimate) the value
of $C_{CW}$. You may, for example, plot $1/k$ vs.\ $T - T_C$ and find the slope of
that line, which is $1/C_{CW}$. You will be using your measured $T_C$ here, which is
easy to identify as the temperature of the sharp cusp in your data.

\notebox{Ramp slowly}{You want to make sure that your temperature ramp is slow enough
to trust your $T_C$ measurement.}

The temperature interval of your data above $T_C$ will be rather narrow, but you will
be able to estimate $C_{CW}$. It will be interesting to also fit the data below the
transition, to see if you can verify the theoretical factor of two difference between
$C_{CW1}$ and $C_{CW2}$ (see Section~\ref{sec:landau}). As we said above, the largest
uncertainties in the value of $C_{CW}$ will come from not knowing the capacitor
geometry well, and from measuring a ceramic sample (although the latter may be in
part remedied by the capacitance correction), rather than from your fits to a
Curie--Weiss law.

Why do we want to find $C_{CW}$? Because there is a theoretical connection between
the value of $C_{CW}$ and microscopic parameters, and we will try to obtain it in the
next section.

\subsection{Derivation of \texorpdfstring{$C_{CW}$}{C-CW} in terms of a microscopic
            dipole moment}\label{sec:ccw}

In a ferromagnetic material microscopic magnetic moments $\mu$ (fundamentally,
spins) are always present; the difference between paramagnetic and ferromagnetic
state is merely in the organization, alignment or misalignment of these pre-existing
moments. In a paraelectric material, a unit cell is not carrying its $p$ in the same
way as an atom carries its spin; in fact the microscopic electric dipole moment $p$
above the $T_C$ should not exist. The lattice symmetry changes at $T_C$, and above
$T_C$ we have a cubic cell with Ti in the center, so that microscopic $p = 0$.

Here we wish to make the following comment: Landau's potential above $T_C$
(Figure~\ref{fig:a5f3}(a)) has only one minimum, and this may suggest that all
electrical dipoles $p$ and magnetic moments $\mu$ must disappear above $T_C$. Yet we
know that individual spins do not disappear at $T > T_C$ in a paramagnetic state.
Thus we realize that the notion of order parameter $m$ (in our case, polarization
$P$) is macroscopic in nature; it does not tell us anything about the existence or
non-existence of microscopic $p$'s above the $T_C$. Our conclusion about
local $p = 0$ in a paraelectric state in \BTO\ is based on the known structure and
symmetry of its unit cell~\cite{ref:jona,ref:blakemore} and not on Landau's theory.

Another difference between magnetic and electric cases is in that quantum mechanics
imposes certain requirements on the quantum numbers and energy levels in the magnetic
case: as is well-known, a degenerate single energy level in $B = 0$ field will split
into two levels for spin one-half electrons in $B \ne 0$
field~\cite{ref:papon,ref:landau1}. This may not be the same in case of $p$'s and
$E$-field: there is no quantum number associated with $p$. Yet we will also use the
same two-level picture here. In other words, we will assume that microscopic dipole
$p$ can point either along the $E$ field or against it, as shown in
Figure~\ref{fig:a5f8}, and will ignore its other possible orientations with respect
to the direction of $E$.

The analysis of the ferromagnetic case can be found in a number of books, e.g.\ in
Refs.~\cite{ref:gittrman,ref:papon,ref:blakemore,ref:ashcroft}. Our derivation here
is similar to these treatments, but somewhat simplified.

Let us therefore consider a material in a ferroelectric state, at $T < T_C$, where we
can assume that one dipole moment $p$ exists in each unit cell of volume $a^3$; thus
in \SI{1}{\cubic\metre} there are $N = 1/a^3$ microscopic dipoles $p$. Let us roughly
estimate the value of a microscopic dipole moment in \BTO. The ``generic'' dipole is
shown in Figure~\ref{fig:dipole}.

\begin{figure}[htbp]
  \centering
  \tikzfig[width=0.3\linewidth]{fig-a5-dipole}{%
    Diagram of an electric dipole: a circle containing minus q sits above a circle
    containing plus q, joined by a straight line. A vertical double-headed arrow
    beside the line marks the separation x, and the dipole moment is given as p
    equals q times x.}
  \caption{A generic electric dipole, $p = qx$.}
  \label{fig:dipole}
\end{figure}

It consists of two opposite electric charges, $-q$ and $+q$, separated by a distance
$x$; the formula is $p = qx$. Whatever ionic displacements happen inside the unit cell
of a ferroelectric, and whatever charge imbalance they will create locally, the total
charge cannot change, and the unit cell must remain neutral. Thus our effective
electrical dipole which represents this polarization, must be neutral as well, with
$-q$ and $+q$ adding to zero. In this formula we treat $q$ as some effective electric
charge, which, we recognize, must be of the order of electron charge $e$, or, at
most, it can be several times larger than $e$. Indeed, we know that in \BTO, Ti ion
carries charge $+4$, and Ba ion carries charge $+2$, while three oxygen ions carry
charge $-6$, with the sum of all these charges being, of course, zero. We will
therefore take $q = 6e$ in our effective dipole.

The second quantity in the dipole formula is $x$, an effective displacement which
represents charge-weighted displacements of all of the charges in the unit cell. In
\BTO, in the first, rough approximation, we can think of $x$ as a displacement of a
Ti\textsuperscript{4+} ion, which in a cubic state resides in the middle of a cell
and shifts up or down (Figure~\ref{fig:a5f8}) in a ferroelectric state. Therefore,
the absolute upper limit on the value of $x$ is $a/2$, half size of the unit
cell. In reality we expect $x$ to be much smaller than this upper limit. The
lower limit is hard to predict; it can be a small fraction of $a$. So we expect to
find $0 < x < a/2$. The length of a cubic unit cell side in \BTO\ is
$a = \SI{0.4}{\nano\metre}$~\cite{ref:jona,ref:blakemore,ref:umich}, which gives us
$0 < x < \SI{0.2}{\nano\metre}$.

Thus we have a rough upper-limit estimate
\begin{equation}
  p_{\max} = q\,x_{\max} = 6e\,x_{\max}
           \approx (\num{9.6e-19})\times(\num{0.2e-9})
           \approx \SI{2e-28}{\coulomb\metre}.
\end{equation}

We also know that in the electric field $E$, potential energy of $p$ is
$U = -(\mathbf{p}\cdot\mathbf{E})$, i.e.\ it changes from the minimum value $(-pE)$
for $p$ being parallel to $E$ to the maximum value $(+pE)$ for $p$ being
anti-parallel to $E$. The thermal energy at room temperature is
$k_BT = \num{1.38e-23}\times 300 = \SI{4.1e-21}{\joule}$. We see that it will take
electric field of \SI{2e7}{\volt\per\metre} to make maximum dipole energy equal to
thermal energy. We conclude that for any reasonable value of electric field $E$,
$U \ll k_BT$.

Now, near the transition, where macroscopic polarization $P$ (the order parameter) is
small, in the absence of an external field $E$, all microscopic $p$'s have the same
(zero) energy. Once the external field $E$ is applied, this energy will split into two
levels, one above and one below the initial level, the energy difference between these
two levels being $\Delta U = pE - (-pE) = 2pE$, as shown in
Figure~\ref{fig:a5levels}.

\begin{figure}[htbp]
  \centering
  \tikzfig[width=0.85\linewidth]{fig-a5-level-splitting}{%
    Energy level diagram. A vertical axis is labeled dipole energy. On the left, a
    single horizontal level carries all N dipoles at zero field. Moving to the right,
    this level splits into two: an upper level, whose population is N over two times
    the exponential of minus delta U over two k-B T, and a lower level, whose
    population is N over two times the exponential of plus delta U over two k-B T.
    Small arrows labeled p on each level show the two dipole orientations. A vertical
    double-headed arrow between the levels is labeled delta U equals 2 p E, and a
    heavy downward arrow at the right marks the applied electric field E.}
  \caption{Splitting of the dipole energy level in an applied electric field.}
  \label{fig:a5levels}
\end{figure}

We still have $\Delta U \ll k_BT$, which means that there will be approximately $N/2$
dipoles on each of the two energy levels. Nonetheless, the small $\Delta U$ is
responsible for the population difference between the two levels: the lower level will
have
\begin{equation}
  N_1 = \frac{N}{2}\exp\!\left(\frac{\Delta U}{2k_BT}\right)
      \approx \frac{N}{2}\left(1 + \frac{\Delta U}{2k_BT}\right)
\end{equation}
dipoles, while the upper level will have
\begin{equation}
  N_2 = \frac{N}{2}\exp\!\left(-\frac{\Delta U}{2k_BT}\right)
      \approx \frac{N}{2}\left(1 - \frac{\Delta U}{2k_BT}\right)
\end{equation}
dipoles, so that $N_1 > N_2$. The total number of dipoles is of course still
$N_1 + N_2 = N$. The relatively small difference $N_1 - N_2$ is what is responsible
for macroscopic polarization $P$. So we have
\begin{equation}
  N_1 - N_2 = \frac{N}{2}\,\frac{\Delta U}{k_BT}
            = \frac{N}{2}\,\frac{2pE}{k_BT}
            = \frac{NpE}{k_BT},
\end{equation}
Macroscopic polarization is
\begin{equation}
  P = p\,(N_1 - N_2) = \frac{N p^{2} E}{k_B T}.
\end{equation}
Finally we will get for susceptibility $\chi_E = P/E$, and for relative dielectric
constant,
\begin{equation}
  \chi_E = \frac{N p^{2}/k_B}{T},
  \qquad
  k = \frac{N p^{2}/\varepsilon_0 k_B}{T}.
\end{equation}

The one difference between the last formula and formula~\eqref{eq:cwk} is in the
temperature dependence: while in~\eqref{eq:cwk} we had $1/(T - T_C)$ (or, below the
transition, $T_C - T$), here we got simply $1/T$. But we did not develop in this
section a full theory of a phase transition; we just assumed that transition takes
place at some temperature; in fact, we could assume that it took place at zero
temperature; then $T - T_C$ would be the same as $T_C$. So we should not be very
surprised that we did not get $1/(T - T_C)$; the correct temperature dependence was
addressed in Landau theory in Section~\ref{sec:landau}. You can also find a
derivation of these formulas with the correct $1/(T - T_C)$ dependence in
Ref.~\cite{ref:papon}.

So we see that we should identify the numerator with $C_{CW}$; by doing so we obtain
\begin{equation}\label{eq:ccw}
  C_{CW} = \frac{N p^{2}}{\varepsilon_0 k_B},
\end{equation}
or, in terms of a microscopic dipole moment $p$,
\begin{equation}\label{eq:pfromccw}
  p = \sqrt{\frac{\varepsilon_0 k_B C_{CW}}{N}},
\end{equation}
where $N = 1/a^3$ and $C_{CW}$ is a measured quantity which can be obtained from
fitting your $k$ vs.\ $T$ data to the Curie--Weiss Law.

\notebox{As an exercise, check the dimensionality of $C_{CW}$ in~\eqref{eq:ccw}}{It
should come out to be K (kelvin), so that $k$ in formula~\eqref{eq:cwk} will indeed
be a \emph{relative} quantity without units.}

\notebox{What you should calculate}{You should find $C_{CW}$ derived from your data
(again, we assume that you applied the capacitance correction), substitute it in
formula~\eqref{eq:pfromccw}, and calculate $p$, and then compare it to our
upper-limit estimate $p_{\max}$ for this quantity. Once you will know $p$, you can
calculate $x$, the effective dipole displacement. This last value can be compared to
the literature: indeed, transition-induced displacements of various atoms in the unit
cell have been studied by X-rays. You should also find in the literature the value of
$C_{CW}$ for single-crystalline \BTO, and, using it in formula~\eqref{eq:pfromccw},
calculate again $p$ and $x$. This time the values will be more realistic, as you will
be using the correct, intrinsic value of $C_{CW}$. Compare the results with your own,
comment on discrepancies.}

\clearpage
% =============================================================================
\begin{thebibliography}{99}
\addcontentsline{toc}{section}{References}
% =============================================================================

\bibitem{ref:lockin}
  \emph{Lock-in amplifier}, Wikipedia.
  \url{http://en.wikipedia.org/wiki/Lock-in_amplifier}
  --- about lock-in amplifiers; at the end of this article you will find
  references to more lock-in tutorials.

\bibitem{ref:accircuits}
  \emph{AC circuits}, School of Physics, University of New South Wales.
  \url{http://www.animations.physics.unsw.edu.au/jw/AC.html}
  --- a tutorial on AC circuits, with animations.

\bibitem{ref:pid}
  \emph{PID controller}, Wikipedia.
  \url{http://en.wikipedia.org/wiki/PID_controller}
  --- about temperature control.

\bibitem{ref:gittrman}
  M.~Gitterman and V.~Halpern,
  \emph{Phase Transitions: A Brief Account with Modern Applications},
  World Scientific (2004).

\bibitem{ref:papon}
  P.~Papon, J.~Leblond and P.~H.~E.~Meijer,
  \emph{The Physics of Phase Transitions}, Springer (2001).

\bibitem{ref:landau1}
  \emph{Landau theory of phase transitions}, lecture notes, University of
  Helsinki.
  \url{http://theory.physics.helsinki.fi/~stafyi/SPI06chap10.pdf}

\bibitem{ref:landau2}
  \emph{Second order phase transitions}, P.~Hadley, TU Graz.
  \url{http://lamp.tu-graz.ac.at/~hadley/ss2/landau/second_order.php}

\bibitem{ref:dobro}
  V.~Dobrosavljevi\'c, \emph{Introduction to Metal--Insulator Transitions}
  (2011), arXiv:1112.6166v1.

\bibitem{ref:mottbook}
  N.~F.~Mott, \emph{Metal--Insulator Transitions},
  Taylor \& Francis Ltd. (1974), ISBN 0-85066-079-3.

\bibitem{ref:mott1949}
  N.~F.~Mott, \emph{Proc. Phys. Soc. (London)} \textbf{A62}, 416 (1949).

\bibitem{ref:jona}
  F.~Jona and G.~Shirane, \emph{Ferroelectric Crystals},
  Pergamon Press / The Macmillan Co. (1962).
  A copy of the relevant pages is provided in the lab.

\bibitem{ref:blakemore}
  J.~S.~Blakemore, \emph{Solid State Physics}, 2nd edition,
  W.~B.~Saunders Co. (1974), Chapter~5, ``Dielectric and magnetic properties of
  solids''.

\bibitem{ref:umich}
  \emph{Phase transitions in ferroelectrics}, Advanced Lab, University of
  Michigan.
  \url{http://instructor.physics.lsa.umich.edu/adv-labs/Ferroelectrics/phase_transitions.pdf}

\bibitem{ref:ashcroft}
  N.~W.~Ashcroft and N.~D.~Mermin, \emph{Solid State Physics},
  Thomson Learning, Toronto (1976).

\bibitem{ref:percolation}
  \emph{Percolation}, Chapter~13 of course notes, Clark University.
  \url{http://physics.clarku.edu/courses/125/gtcdraft/chap13.pdf}

\bibitem{ref:bergland}
  C.~N.~Berglund and H.~J.~Guggenheim,
  ``Electronic properties of \VOtwo\ near the semiconductor--metal transition'',
  \emph{Phys. Rev.} \textbf{185}(3), 1022--1033 (1969).

\bibitem{ref:gurvitch}
  M.~Gurvitch, S.~Luryi, A.~Polyakov and A.~Shabalov,
  ``Nonhysteretic phenomena in the metal--semiconductor phase-transition loop of
  \VOtwo\ films for bolometric sensor applications'',
  \emph{IEEE Transactions on Nanotechnology} \textbf{9}, 647--652 (2010).

\bibitem{ref:lanskaya}
  For example, see T.~G.~Lanskaya, I.~A.~Merkulov and F.~A.~Chudnovskii,
  ``Hysteresis effects at the semiconductor--metal phase transition in vanadium
  oxides'', \emph{Sov. Phys. Solid State} \textbf{20}(2), 193 (1978), and a
  number of more recent papers treating hysteresis in \VOtwo.

\bibitem{ref:doitpoms}
  \emph{Ferroelectric materials: phase changes}, DoITPoMS, University of
  Cambridge.
  \url{http://www.doitpoms.ac.uk/tlplib/ferroelectrics/phase_changes.php}
  --- about the structure of \BTO, with very nice graphics and illustrations.

\end{thebibliography}

\bigskip
\noindent\rule{\linewidth}{0.4pt}

\medskip
\noindent\small
Sections relating to ferroelectricity in \BTO\ were created by Laszlo Mihaly and
in part by Michael Gurvitch; sections relating to phase-transition theory and the
metal--insulator transition in \VOtwo\ were created by Michael Gurvitch, both in
2012. Thanks to Xu Du, Matthew Dawber, Alexander Shabalov and Carlos Marques for
contributions to the design and realization of the measurement and for useful
comments. The 2026 edition of this manual was typeset into \LaTeX \ with Claude
Opus 4.5 and updated by Logan Levack.

\end{document}
